\documentclass[12pt]{article}
% =========================================================
% UNIDAD 2: FUNCIONES
% =========================================================
\usepackage[utf8]{inputenc}
\usepackage[T1]{fontenc}
\usepackage[spanish,mexico]{babel}
\usepackage{amsmath,amssymb,amsthm,mathtools}
\usepackage{geometry}
\usepackage{enumitem}
\usepackage{booktabs}
\usepackage{longtable}
\usepackage{array}
\usepackage{xcolor}
\usepackage{hyperref}
\usepackage{graphicx}
\usepackage{tikz}
\usepackage{pgfplots}
\usepackage{multicol}
\usepackage{fancyhdr}
% =========================================================
\pgfplotsset{compat=1.18}
\pgfplotsset{unbounded coords=jump}
\geometry{letterpaper,margin=2.4cm}
\setlength{\parindent}{0pt}
\setlength{\parskip}{0.55em}
\setlist[itemize]{leftmargin=1.7em}
\setlist[enumerate]{leftmargin=2em}
% =========================================================
\definecolor{azul}{RGB}{25,70,125}
\hypersetup{
    colorlinks=true,
    linkcolor=azul,
    urlcolor=azul,
    citecolor=azul
}
% =========================================================
\pagestyle{fancy}
\fancyhf{}
\fancyhead[L]{Tecnológico Nacional de México}
\fancyhead[R]{Cálculo Diferencial}
\fancyfoot[L]{Carlos Ernesto Martínez Rodríguez}
\fancyfoot[C]{\thepage}
\fancyfoot[R]{Unidad 2: Funciones}
\renewcommand{\headrulewidth}{0.4pt}
\renewcommand{\footrulewidth}{0.4pt}
% =========================================================
\newtheorem{definicion}{Definición}[section]
\newtheorem{teorema}{Teorema}[section]
\newtheorem{proposicion}{Proposición}[section]
\newtheorem{corolario}{Corolario}[section]
\theoremstyle{remark}
\newtheorem{observacion}{Observación}[section]
\newtheorem{ejemplo}{Ejemplo}[section]
\newtheorem{nota}{Nota}[section]
% =========================================================
\newcommand{\R}{\mathbb{R}}
\newcommand{\Dom}{\operatorname{Dom}}
\newcommand{\Ima}{\operatorname{Im}}
\newcommand{\Ran}{\operatorname{Ran}}
\newcommand{\id}{\operatorname{id}}
% =========================================================
\title{\textbf{Notas para el curso de Cálculo Diferencial}\\[0.35cm]
\Large Unidad 2: Funciones}
\author{Carlos Ernesto Martínez Rodríguez}
\date{}
% =========================================================
\begin{document}

\maketitle
\tableofcontents
\newpage

% =========================================================
\section*{Propósito de la unidad}
\addcontentsline{toc}{section}{Propósito de la unidad}

El concepto de función constituye uno de los lenguajes fundamentales del Cálculo. Los límites, la continuidad, la derivación y sus aplicaciones se formulan sobre funciones; por ello, antes de estudiar dichos conceptos es necesario comprender con precisión qué es una función, cuáles son los conjuntos que intervienen en su definición, cómo se representa, cómo se opera con ella y cómo puede emplearse para construir modelos matemáticos. El objetivo de esta unidad es analizar la definición de función real, identificar distintos tipos de funciones y sus representaciones gráficas.

% =========================================================
\section{Relaciones y funciones}
% =========================================================

%<<>><<>><<>><<>><<>><<>><<>><<>><<>><<>><<>><<>>
\subsection*{Producto cartesiano}
%<<>><<>><<>><<>><<>><<>><<>><<>><<>><<>><<>><<>>

\begin{definicion}[Producto cartesiano]
Sean $A$ y $B$ dos conjuntos. El \emph{producto cartesiano} de $A$ y $B$ es
\[
A\times B=\{(a,b):a\in A,\ b\in B\}.
\]
Sus elementos son pares ordenados.
\end{definicion}

El calificativo ordenado es importante, pues en general $(a,b)\neq(b,a)$.
Dos pares ordenados son iguales si y sólo si coinciden sus componentes:
\[
(a,b)=(c,d)
\quad\Longleftrightarrow\quad
a=c\ \text{y}\ b=d.
\]

\begin{ejemplo}
Si
\[
A=\{1,2\},\qquad B=\{a,b,c\},
\]
entonces
\[
A\times B=
\{(1,a),(1,b),(1,c),(2,a),(2,b),(2,c)\}.
\]
Por tanto, $|A\times B|=6=|A||B|$.
\end{ejemplo}
%<<>><<>><<>><<>><<>><<>><<>><<>><<>><<>><<>><<>>
\subsection*{Relaciones}
%<<>><<>><<>><<>><<>><<>><<>><<>><<>><<>><<>><<>>
\begin{definicion}[Relación]
Una relación de $A$ en $B$ es cualquier subconjunto $R\subseteq A\times B$.
\end{definicion}

Una relación permite asociar elementos de $A$ con elementos de $B$. No se
exige que todos los elementos de $A$ estén relacionados, ni que un elemento de $A$ tenga una única pareja.

\begin{ejemplo}
Sean
\[
A=\{1,2,3\},\qquad B=\{a,b,c\}.
\]
El conjunto
\[
R=\{(1,a),(1,b),(3,c)\}
\]
es una relación de $A$ en $B$. Observe que $1$ está relacionado con dos
elementos de $B$ y que $2$ no está relacionado con ninguno.
\end{ejemplo}

%<<>><<>><<>><<>><<>><<>><<>><<>><<>><<>><<>><<>>
\subsection{Definición formal de función}
%<<>><<>><<>><<>><<>><<>><<>><<>><<>><<>><<>><<>>
\begin{definicion}[Función]
Sean $A$ y $B$ conjuntos no vacíos. Una \emph{función} $f$ de $A$ en $B$,
denotada por
\[
f:A\longrightarrow B,
\]
es una relación $f\subseteq A\times B$ que satisface:

\begin{itemize}
    \item para todo $x\in A$ existe algún $y\in B$ tal que $(x,y)\in f$;
    \item el elemento $y$ es único.
\end{itemize}

Equivalentemente,
\[
\forall x\in A\;\exists!\,y\in B
\quad\text{tal que}\quad y=f(x).
\]
\end{definicion}

El símbolo $\exists!$ significa \textit{existe uno y sólo uno}. Así, una función satisface dos requisitos:

\begin{itemize}
    \item \textbf{existencia:} todo elemento del dominio debe tener imagen;
    \item \textbf{unicidad:} cada elemento del dominio debe tener una sola imagen.
\end{itemize}

Es importante notar que distintos elementos del dominio sí pueden tener la misma imagen. Lo prohibido es que un mismo elemento del dominio tenga dos imágenes diferentes.

\begin{ejemplo}
La correspondencia
\[
f:\R\to\R,\qquad f(x)=x^2
\]
es una función. Por ejemplo,
\[
f(2)=4,\qquad f(-2)=4.
\]
Que dos elementos diferentes tengan la misma imagen no contradice la
definición.

\textbf{Interpretación gráfica.} Los puntos $x=2$ y $x=-2$ producen la misma
altura $y=4$. Esto está permitido: la prueba de la recta vertical sigue
cumpliéndose porque para cada valor de $x$ existe una sola altura.
\begin{center}
\begin{tikzpicture}
\begin{axis}[width=.72\textwidth,height=5.4cm,axis lines=middle,
 xmin=-3.5,xmax=3.5,ymin=-1,ymax=10,xlabel={$x$},ylabel={$y$},grid=both,samples=160]
\addplot[thick,domain=-3.2:3.2] {x^2};
\addplot[only marks,mark=*] coordinates {(-2,4) (2,4)};
\draw[dashed] (axis cs:-2,4)--(axis cs:2,4);
\node[above left] at (axis cs:-2,4) {$(-2,4)$};
\node[above right] at (axis cs:2,4) {$(2,4)$};
\end{axis}
\end{tikzpicture}
\end{center}
\end{ejemplo}

\begin{ejemplo}[Una relación que no es función]
Considere
\[
R=\{(1,2),(1,3),(2,5)\}.
\]
No puede ser una función cuyo dominio contenga a $1$, pues el mismo elemento $1$ estaría asociado con dos valores diferentes:
\[
1\mapsto2,\qquad 1\mapsto3.
\]
La unicidad falla.

\textbf{Representación en el plano.} Los pares de la relación pueden verse como
puntos. La recta vertical $x=1$ pasa por dos de ellos, $(1,2)$ y $(1,3)$;
por eso esta relación no puede ser la gráfica de una función de $x$.
\begin{center}
\begin{tikzpicture}
\begin{axis}[width=.62\textwidth,height=5cm,axis lines=middle,
 xmin=0,xmax=3,ymin=0,ymax=6,xtick={1,2},ytick={2,3,5},xlabel={$x$},ylabel={$y$},grid=both]
\addplot[only marks,mark=*,mark size=2.5pt] coordinates {(1,2) (1,3) (2,5)};
\addplot[dashed,domain=0:6] coordinates {(1,0) (1,6)};
\end{axis}
\end{tikzpicture}
\end{center}
\end{ejemplo}
%<<>><<>><<>><<>><<>><<>><<>><<>><<>><<>><<>><<>>
\subsection{Dominio, codominio e imagen}
%<<>><<>><<>><<>><<>><<>><<>><<>><<>><<>><<>><<>>
\begin{definicion}
Sea $f:A\to B$. Entonces:

\begin{itemize}
    \item el conjunto $A$ se llama \emph{dominio} de $f$;
    \item el conjunto $B$ se llama \emph{codominio} de $f$;
    \item si $x\in A$, el elemento $f(x)\in B$ se llama \emph{imagen} de $x$;
    \item el conjunto $f(A)=\{f(x):x\in A\}$ se llama \emph{imagen}, \emph{recorrido} o \emph{rango} de $f$.
\end{itemize}
\end{definicion}

Siempre se cumple $f(A)\subseteq B$, pero no necesariamente $f(A)=B$.

\begin{ejemplo}
Sea
\[
f:\R\to\R,\qquad f(x)=x^2.
\]
Su dominio es $\R$, su codominio es $\R$, pero su imagen es
\[
f(\R)=[0,\infty).
\]
Por tanto,
\[
f(\R)\neq\R.
\]

\textbf{Interpretación gráfica.} La parábola nunca toma valores negativos.
La gráfica permite distinguir visualmente el codominio $\R$ de la imagen
realmente alcanzada, $[0,\infty)$.
\begin{center}
\begin{tikzpicture}
\begin{axis}[width=.70\textwidth,height=5.3cm,axis lines=middle,
 xmin=-3.5,xmax=3.5,ymin=-3,ymax=10,xlabel={$x$},ylabel={$y$},grid=both,samples=160]
\addplot[thick,domain=-3.2:3.2] {x^2};
\addplot[only marks,mark=*] coordinates {(0,0)};
\end{axis}
\end{tikzpicture}
\end{center}
\end{ejemplo}

\begin{observacion}
Dominio, codominio e imagen son conceptos diferentes. El codominio forma parte de la definición de la función; la imagen es el conjunto de valores que la función realmente alcanza.
\end{observacion}
%<<>><<>><<>><<>><<>><<>><<>><<>><<>><<>><<>><<>>
\subsection{Imagen inversa}
%<<>><<>><<>><<>><<>><<>><<>><<>><<>><<>><<>><<>>
\begin{definicion}[Imagen inversa de un conjunto]
Sea $f:A\to B$ y sea $C\subseteq B$. La Imagen inversa (preimagen) de $C$ bajo $f$ se define por
\[
f^{-1}(C)=\{x\in A:f(x)\in C\}.
\]
\end{definicion}

Aquí el símbolo $f^{-1}(C)$ no presupone que exista una función inversa.

\begin{ejemplo}
Sea $f(x)=x^2$ y $C=[1,4]$. Entonces
\[
f^{-1}(C)=\{x\in\R:1\leq x^2\leq4\}
=[-2,-1]\cup[1,2].
\]

\textbf{Interpretación gráfica.} La condición $1\leq f(x)\leq4$ selecciona
los puntos de la parábola comprendidos entre las alturas $y=1$ y $y=4$.
Al proyectarlos sobre el eje $x$ aparecen exactamente los dos intervalos
$[-2,-1]$ y $[1,2]$.
\begin{center}
\begin{tikzpicture}
\begin{axis}[width=.72\textwidth,height=5.5cm,axis lines=middle,
 xmin=-3,xmax=3,ymin=-1,ymax=6,xlabel={$x$},ylabel={$y$},grid=both,samples=180]
\addplot[thick,domain=-2.5:2.5] {x^2};
\addplot[dashed,domain=-3:3] {1};
\addplot[dashed,domain=-3:3] {4};
\addplot[very thick,domain=-2:-1] {x^2};
\addplot[very thick,domain=1:2] {x^2};
\addplot[only marks,mark=*] coordinates {(-2,4) (-1,1) (1,1) (2,4)};
\end{axis}
\end{tikzpicture}
\end{center}
\end{ejemplo}

% =========================================================
\section{Funciones reales de variable real}
% =========================================================

En Cálculo Diferencial trabajaremos principalmente con funciones cuyos
argumentos y valores son números reales.

\begin{definicion}
Una función real de variable real es una función
\[
f:A\to\R,
\]
donde $A\subseteq\R$.
\end{definicion}

Cuando una función se presenta solamente mediante una expresión algebraica y
no se especifica explícitamente su dominio, adoptaremos como dominio el
\emph{dominio natural}: el conjunto más grande de números reales para los que
la expresión tiene sentido.

% =========================================================
\section{Evaluación de funciones}

Evaluar una función consiste en determinar la imagen de un elemento particular
del dominio.

\begin{ejemplo}
Sea
\[
f(x)=3x^2-2x+5.
\]
Entonces
\[
f(2)=3(2)^2-2(2)+5=12-4+5=13.
\]
Asimismo,
\[
f(-1)=3(-1)^2-2(-1)+5=3+2+5=10.
\]

\textbf{Interpretación gráfica.} Evaluar $f(c)$ significa localizar la altura
de la gráfica sobre $x=c$. En este ejemplo, los cálculos anteriores producen
los puntos $(2,13)$ y $(-1,10)$.
\begin{center}
\begin{tikzpicture}
\begin{axis}[width=.72\textwidth,height=5.4cm,axis lines=middle,
 xmin=-2.5,xmax=3.5,ymin=0,ymax=24,xlabel={$x$},ylabel={$y$},grid=both,samples=160]
\addplot[thick,domain=-2.2:3.2] {3*x^2-2*x+5};
\addplot[only marks,mark=*] coordinates {(2,13) (-1,10)};
\node[above right] at (axis cs:2,13) {$(2,13)$};
\node[above left] at (axis cs:-1,10) {$(-1,10)$};
\end{axis}
\end{tikzpicture}
\end{center}
\end{ejemplo}

La sustitución debe hacerse sobre \emph{todas} las apariciones de la variable.

\begin{ejemplo}
Para
\[
f(x)=2x^2-3x+1,
\]
se tiene
\[
f(a)=2a^2-3a+1,
\]
y
\begin{align*}
f(a+h)
 &=2(a+h)^2-3(a+h)+1\\
 &=2(a^2+2ah+h^2)-3a-3h+1\\
 &=2a^2+4ah+2h^2-3a-3h+1.
\end{align*}
\end{ejemplo}

La expresión $f(x+h)$ aparecerá de manera fundamental en la definición de
derivada. Conviene distinguirla de $f(x)+h$.

\begin{ejemplo}
Si $f(x)=x^2$, entonces
\[
f(x+h)=(x+h)^2=x^2+2xh+h^2,
\]
mientras que
\[
f(x)+h=x^2+h.
\]
En general,
\[
f(x+h)\neq f(x)+h.
\]
\end{ejemplo}

\subsection{Cociente incremental como preparación para la derivada}

Aunque la derivada se estudiará posteriormente, resulta útil practicar desde
ahora la simplificación
\[
\frac{f(x+h)-f(x)}{h},\qquad h\neq0.
\]

\begin{ejemplo}
Sea $f(x)=x^2+3x$. Entonces
\[
f(x+h)=(x+h)^2+3(x+h)
=x^2+2xh+h^2+3x+3h.
\]
Por tanto,
\begin{align*}
f(x+h)-f(x)
&=x^2+2xh+h^2+3x+3h-(x^2+3x)\\
&=2xh+h^2+3h\\
&=h(2x+h+3).
\end{align*}
Para $h\neq0$,
\[
\frac{f(x+h)-f(x)}{h}=2x+h+3.
\]
\end{ejemplo}


% =========================================================
\section{Definición previa de las funciones logarítmica y exponencial}
% =========================================================

Antes de utilizar funciones logarítmicas para determinar dominios, conviene
precisar qué entendemos por logaritmo y por función exponencial.

\subsection{Función logarítmica}

Fijemos una base
\[
a>0,\qquad a\neq1.
\]
Suponiendo conocida la potenciación real estudiada en cursos de álgebra,
definimos el logaritmo en base $a$ mediante la equivalencia
\[
\boxed{\log_a y=x\quad\Longleftrightarrow\quad a^x=y},
\qquad y>0.
\]

Esta equivalencia debe leerse en ambos sentidos. Decir
\[
\log_a y=x
\]
significa exactamente que $x$ es el exponente al que debemos elevar la base
$a$ para obtener $y$.

Por ejemplo,
\[
\log_2 8=3
\quad\Longleftrightarrow\quad
2^3=8,
\]
\[
\log_{10}1000=3
\quad\Longleftrightarrow\quad
10^3=1000,
\]
y
\[
\log_3\!\left(\frac19\right)=-2
\quad\Longleftrightarrow\quad
3^{-2}=\frac19.
\]

Como $a^x>0$ para todo $x\in\R$, un logaritmo real solamente puede recibir
argumentos positivos. Por ello,
\[
\boxed{\Dom(\log_a)=(0,\infty)},
\qquad
\boxed{\Ima(\log_a)=\R}.
\]

Si $a>1$, $\log_a x$ es estrictamente creciente; si $0<a<1$, es estrictamente
decreciente. En ambos casos pasa por $(1,0)$ porque $a^0=1$.

\subsection{La función exponencial como inversa del logaritmo}

La equivalencia anterior permite presentar la función exponencial como la
función inversa de $\log_a$. Definimos
\[
\boxed{\exp_a=(\log_a)^{-1}}.
\]
Es decir,
\[
\boxed{\exp_a(x)=a^x}.
\]
Por tanto,
\[
\boxed{y=a^x\quad\Longleftrightarrow\quad x=\log_a y}.
\]

De la inversión se intercambian dominio e imagen:
\[
\Dom(\exp_a)=\R,
\qquad
\Ima(\exp_a)=(0,\infty).
\]
Además,
\[
\log_a(a^x)=x
\]
y, para $y>0$,
\[
a^{\log_a y}=y.
\]

Geométricamente, las gráficas de una función y de su inversa son reflejos
respecto de la recta $y=x$. Para $a=2$:

\begin{center}
\begin{tikzpicture}
\begin{axis}[
 width=.78\textwidth,height=6.2cm,axis lines=middle,
 xmin=-4,xmax=5,ymin=-4,ymax=8,xlabel={$x$},ylabel={$y$},
 grid=both,samples=180,legend style={font=\small,at={(0.02,0.98)},anchor=north west}]
\addplot[thick,domain=-4:3] {2^x};
\addlegendentry{$y=2^x$}
\addplot[thick,domain=0.07:5] {ln(x)/ln(2)};
\addlegendentry{$y=\log_2x$}
\addplot[dashed,domain=-4:5] {x};
\addlegendentry{$y=x$}
\end{axis}
\end{tikzpicture}
\end{center}

\begin{observacion}
La definición
\[
\log_a y=x\Longleftrightarrow a^x=y
\]
es la herramienta que utilizaremos para pasar de forma logarítmica a forma
exponencial y viceversa. También explica inmediatamente por qué el argumento
de un logaritmo real debe ser estrictamente positivo.
\end{observacion}

% =========================================================
\section{Determinación del dominio}

\subsection{Funciones polinomiales}

Todo polinomio
\[
p(x)=a_nx^n+\cdots+a_1x+a_0
\]
está definido para todo número real. Por tanto,
\[
\Dom(p)=\R.
\]

\subsection{Funciones racionales}

Una función racional tiene la forma
\[
f(x)=\frac{p(x)}{q(x)},
\]
donde $p$ y $q$ son polinomios y $q$ no es el polinomio cero.

Como la división entre cero no está definida,
\[
\Dom(f)=\{x\in\R:q(x)\neq0\}.
\]

\begin{ejemplo}
Determine el dominio de
\[
f(x)=\frac{x+1}{x^2-5x+6}.
\]
Factorizando,
\[
x^2-5x+6=(x-2)(x-3).
\]
El denominador se anula para $x=2$ y $x=3$. Por tanto,
\[
\Dom(f)=\R\setminus\{2,3\}
=(-\infty,2)\cup(2,3)\cup(3,\infty).
\]

\textbf{Interpretación gráfica.} Las rectas $x=2$ y $x=3$ son asíntotas
verticales de la función; precisamente esos valores quedan fuera del dominio.
\begin{center}
\begin{tikzpicture}
\begin{axis}[width=.72\textwidth,height=5.6cm,axis lines=middle,
 xmin=-2,xmax=7,ymin=-6,ymax=6,xlabel={$x$},ylabel={$y$},grid=both,samples=220]
\addplot[thick,domain=-2:1.95] {(x+1)/((x-2)*(x-3))};
\addplot[thick,domain=2.05:2.95] {(x+1)/((x-2)*(x-3))};
\addplot[thick,domain=3.05:7] {(x+1)/((x-2)*(x-3))};
\addplot[dashed] coordinates {(2,-6) (2,6)};
\addplot[dashed] coordinates {(3,-6) (3,6)};
\end{axis}
\end{tikzpicture}
\end{center}
\end{ejemplo}

\subsection{Radicales de índice par}

Para que $\sqrt{g(x)}$ sea real se requiere
\[
g(x)\geq0.
\]

\begin{ejemplo}
Determine el dominio de
\[
f(x)=\sqrt{5-2x}.
\]
Se exige
\[
5-2x\geq0.
\]
Entonces
\[
-2x\geq-5
\]
y al dividir entre $-2$ se invierte la desigualdad:
\[
x\leq\frac52.
\]
Así,
\[
\Dom(f)=\left(-\infty,\frac52\right].
\]

\textbf{Interpretación gráfica.} La gráfica termina en $x=5/2$; a la derecha
no existen valores reales de $\sqrt{5-2x}$.
\begin{center}
\begin{tikzpicture}
\begin{axis}[width=.72\textwidth,height=5.4cm,axis lines=middle,
 xmin=-4,xmax=4,ymin=-1,ymax=4.5,xlabel={$x$},ylabel={$y$},grid=both,samples=180]
\addplot[thick,domain=-4:2.5] {sqrt(5-2*x)};
\addplot[only marks,mark=*] coordinates {(2.5,0)};
\end{axis}
\end{tikzpicture}
\end{center}
\end{ejemplo}

\subsection{Funciones logarítmicas}

Para el logaritmo real se requiere argumento estrictamente positivo:
\[
\log_a(g(x))
\quad\Longrightarrow\quad
g(x)>0,
\]
donde $a>0$ y $a\neq1$.

\begin{ejemplo}
Determine el dominio de
\[
f(x)=\ln(x^2-4).
\]
Debemos resolver
\[
x^2-4>0.
\]
Factorizando,
\[
(x-2)(x+2)>0,
\]
por lo que
\[
x<-2\quad\text{o}\quad x>2.
\]
Así,
\[
\Dom(f)=(-\infty,-2)\cup(2,\infty).
\]

La desigualdad puede estudiarse sistemáticamente mediante una tabla de signos:
\[
\begin{array}{c|ccccc}
x&(-\infty,-2)&-2&(-2,2)&2&(2,\infty)\\\\ \hline
x-2&-&-&-&0&+\\\\
x+2&-&0&+&+&+\\\\ \hline
(x-2)(x+2)&+&0&-&0&+
\end{array}
\]
Como el argumento del logaritmo debe satisfacer $x^2-4>0$, los ceros
$x=\pm2$ no se incluyen.

\textbf{Interpretación gráfica.} La función tiene dos ramas, únicamente donde
$x^2-4$ es positivo, y presenta asíntotas verticales en $x=-2$ y $x=2$.
\begin{center}
\begin{tikzpicture}
\begin{axis}[width=.72\textwidth,height=5.5cm,axis lines=middle,
 xmin=-6,xmax=6,ymin=-5,ymax=4,xlabel={$x$},ylabel={$y$},grid=both,samples=220]
\addplot[thick,domain=-6:-2.02] {ln(x^2-4)};
\addplot[thick,domain=2.02:6] {ln(x^2-4)};
\addplot[dashed] coordinates {(-2,-5) (-2,4)};
\addplot[dashed] coordinates {(2,-5) (2,4)};
\end{axis}
\end{tikzpicture}
\end{center}
\end{ejemplo}

\subsection{Restricciones simultáneas}

Cuando una expresión contiene varias operaciones restrictivas, todas las
condiciones deben cumplirse simultáneamente.

\begin{ejemplo}
Determine el dominio de
\[
f(x)=\frac{\sqrt{x-1}}{x-3}.
\]
Se requiere
\[
x-1\geq0
\quad\text{y}\quad
x-3\neq0.
\]
La primera condición produce $x\geq1$ y la segunda excluye $x=3$. Por tanto,
\[
\Dom(f)=[1,3)\cup(3,\infty).
\]

\textbf{Interpretación gráfica.} El punto inicial $(1,0)$ sí pertenece a la
gráfica, mientras que $x=3$ queda excluido y produce una asíntota vertical.
\begin{center}
\begin{tikzpicture}
\begin{axis}[width=.72\textwidth,height=5.5cm,axis lines=middle,
 xmin=0,xmax=8,ymin=-4,ymax=4,xlabel={$x$},ylabel={$y$},grid=both,samples=220]
\addplot[thick,domain=1:2.96] {sqrt(x-1)/(x-3)};
\addplot[thick,domain=3.04:8] {sqrt(x-1)/(x-3)};
\addplot[dashed] coordinates {(3,-4) (3,4)};
\addplot[only marks,mark=*] coordinates {(1,0)};
\end{axis}
\end{tikzpicture}
\end{center}
\end{ejemplo}

\begin{ejemplo}
Determine el dominio de
\[
f(x)=\sqrt{\frac{x-1}{x+2}}.
\]
Se requiere
\[
\frac{x-1}{x+2}\geq0,
\qquad x\neq-2.
\]
Los puntos críticos son $-2$ y $1$. Mediante una tabla de signos:
\[
\frac{x-1}{x+2}\geq0
\quad\Longleftrightarrow\quad
x\in(-\infty,-2)\cup[1,\infty).
\]
Por tanto,
\[
\Dom(f)=(-\infty,-2)\cup[1,\infty).
\]

Aquí la frase ``mediante una tabla de signos'' significa construirla
explícitamente. Los puntos críticos son $-2$ (anula el denominador) y $1$
(anula el numerador):
\[
\begin{array}{c|ccccc}
x&(-\infty,-2)&-2&(-2,1)&1&(1,\infty)\\\\ \hline
x-1&-&-&-&0&+\\\\
x+2&-&0&+&+&+\\\\ \hline
\dfrac{x-1}{x+2}&+&\text{N.D.}&-&0&+
\end{array}
\]
Para la raíz cuadrada necesitamos cociente $\ge0$: elegimos los intervalos
con signo positivo y también $x=1$, donde el cociente vale cero; nunca podemos
incluir $x=-2$.

\textbf{Interpretación gráfica.}
\begin{center}
\begin{tikzpicture}
\begin{axis}[width=.72\textwidth,height=5.5cm,axis lines=middle,
 xmin=-8,xmax=7,ymin=-1,ymax=4,xlabel={$x$},ylabel={$y$},grid=both,samples=220]
\addplot[thick,domain=-8:-2.03] {sqrt((x-1)/(x+2))};
\addplot[thick,domain=1:7] {sqrt((x-1)/(x+2))};
\addplot[dashed] coordinates {(-2,-1) (-2,4)};
\addplot[only marks,mark=*] coordinates {(1,0)};
\end{axis}
\end{tikzpicture}
\end{center}
\end{ejemplo}

% =========================================================
\section{Gráfica de una función}

\begin{definicion}[Gráfica]
Sea $f:A\to\R$, con $A\subseteq\R$. La gráfica de $f$ es el conjunto
\[
G_f=\{(x,f(x)):x\in A\}\subseteq\R^2.
\]
\end{definicion}

La gráfica no es simplemente un dibujo: matemáticamente es un conjunto de
pares ordenados.

\subsection{Prueba de la recta vertical}

Una curva del plano representa la gráfica de una función $y=f(x)$ si toda
recta vertical $x=a$ intersecta la curva en, a lo más, un punto.

La justificación proviene directamente de la unicidad de la definición de
función: para cada $x$ puede existir como máximo un valor $y=f(x)$.

\subsection{Intersecciones con los ejes}

Para encontrar la intersección con el eje $y$, si $0$ pertenece al dominio, se
calcula
\[
f(0).
\]

Las intersecciones con el eje $x$ son las soluciones de
\[
f(x)=0.
\]

\begin{ejemplo}
Sea
\[
f(x)=x^2-5x+6.
\]
La intersección con el eje $y$ es
\[
f(0)=6,
\]
por lo que se obtiene $(0,6)$.

Para el eje $x$:
\[
x^2-5x+6=0,
\]
\[
(x-2)(x-3)=0.
\]
Así, las intersecciones son
\[
(2,0),\qquad(3,0).
\]

La gráfica permite comprobar simultáneamente los tres puntos obtenidos:
\begin{center}
\begin{tikzpicture}
\begin{axis}[width=.72\textwidth,height=5.5cm,axis lines=middle,
 xmin=-1,xmax=6,ymin=-2,ymax=8,xlabel={$x$},ylabel={$y$},grid=both,samples=160]
\addplot[thick,domain=-1:6] {x^2-5*x+6};
\addplot[only marks,mark=*] coordinates {(0,6) (2,0) (3,0)};
\end{axis}
\end{tikzpicture}
\end{center}
\end{ejemplo}

% =========================================================
\section{Funciones pares e impares}

\begin{definicion}[Función par]
Una función $f$ cuyo dominio es simétrico respecto de $0$ es \emph{par} si
\[
f(-x)=f(x)
\]
para todo $x$ de su dominio.
\end{definicion}

Geométricamente, una función par tiene simetría respecto del eje $y$.

\begin{definicion}[Función impar]
Una función $f$ cuyo dominio es simétrico respecto de $0$ es \emph{impar} si
\[
f(-x)=-f(x)
\]
para todo $x$ de su dominio.
\end{definicion}

Geométricamente, una función impar tiene simetría respecto del origen.

\begin{ejemplo}
Para
\[
f(x)=x^4+2x^2,
\]
\[
f(-x)=(-x)^4+2(-x)^2=x^4+2x^2=f(x).
\]
Por tanto, $f$ es par.

Gráficamente se observa la simetría respecto del eje $y$:
\begin{center}
\begin{tikzpicture}
\begin{axis}[width=.68\textwidth,height=5.2cm,axis lines=middle,
 xmin=-2,xmax=2,ymin=-1,ymax=18,xlabel={$x$},ylabel={$y$},grid=both,samples=150]
\addplot[thick,domain=-2:2] {x^4+2*x^2};
\end{axis}
\end{tikzpicture}
\end{center}
\end{ejemplo}

\begin{ejemplo}
Para
\[
g(x)=x^3-4x,
\]
\[
g(-x)=(-x)^3-4(-x)=-x^3+4x=-g(x).
\]
Por tanto, $g$ es impar.

La gráfica presenta simetría respecto del origen:
\begin{center}
\begin{tikzpicture}
\begin{axis}[width=.68\textwidth,height=5.2cm,axis lines=middle,
 xmin=-3,xmax=3,ymin=-8,ymax=8,xlabel={$x$},ylabel={$y$},grid=both,samples=170]
\addplot[thick,domain=-3:3] {x^3-4*x};
\end{axis}
\end{tikzpicture}
\end{center}
\end{ejemplo}

\begin{observacion}
Una función no tiene por qué ser par o impar. Por ejemplo,
\[
h(x)=x^2+x
\]
no satisface ninguna de las dos identidades.
\end{observacion}

% =========================================================
\section{Crecimiento y decrecimiento}

\begin{definicion}
Sea $f:A\to\R$ y sea $I\subseteq A$ un intervalo.

\begin{enumerate}
    \item $f$ es \emph{creciente} en $I$ si
    \[
    x_1<x_2\quad\Longrightarrow\quad f(x_1)\leq f(x_2).
    \]
    \item $f$ es \emph{estrictamente creciente} si
    \[
    x_1<x_2\quad\Longrightarrow\quad f(x_1)<f(x_2).
    \]
    \item $f$ es \emph{decreciente} si
    \[
    x_1<x_2\quad\Longrightarrow\quad f(x_1)\geq f(x_2).
    \]
\end{enumerate}
\end{definicion}

En unidades posteriores la derivada proporcionará una herramienta sistemática
para determinar intervalos de crecimiento y decrecimiento. En esta etapa se
trabajará con argumentos algebraicos y gráficos.

% =========================================================
\section{Funciones elementales}

\subsection{Función constante}

\[
f(x)=c,\qquad c\in\R.
\]
Su gráfica es una recta horizontal y
\[
\Ima(f)=\{c\}.
\]

\subsection{Función identidad}

\[
\id_{\R}(x)=x.
\]
Su gráfica es la recta $y=x$.

\subsection{Función afín}

\[
f(x)=mx+b.
\]
El parámetro $m$ determina la pendiente y $b$ la intersección con el eje
vertical.

Si $m>0$, la función es estrictamente creciente; si $m<0$, es estrictamente
decreciente; si $m=0$, es constante.

\subsection{Función cuadrática}

Una función cuadrática tiene la forma
\[
f(x)=ax^2+bx+c,\qquad a\neq0.
\]
Su gráfica es una parábola.

Completando el cuadrado:
\[
f(x)
=a\left(x+\frac{b}{2a}\right)^2
+\left(c-\frac{b^2}{4a}\right).
\]
El vértice tiene coordenadas
\[
\left(
-\frac{b}{2a},
f\left(-\frac{b}{2a}\right)
\right).
\]

\subsection{Funciones polinomiales}

\[
p(x)=a_nx^n+a_{n-1}x^{n-1}+\cdots+a_1x+a_0,
\qquad a_n\neq0.
\]
El entero $n$ es el grado del polinomio.

\subsection{Funciones racionales}

\[
f(x)=\frac{p(x)}{q(x)},\qquad q(x)\neq0.
\]

\subsection{Función valor absoluto}

\[
f(x)=|x|=
\begin{cases}
x,&x\geq0,\\
-x,&x<0.
\end{cases}
\]

\subsection{Funciones potencia y radicales}

Las funciones
\[
f(x)=x^\alpha
\]
presentan distintos dominios y comportamientos según el exponente $\alpha$.

En particular,
\[
f(x)=\sqrt{x}=x^{1/2}
\]
tiene dominio $[0,\infty)$.

\subsection{Funciones exponenciales y logarítmicas}

Como se definió previamente, para $a>0$, $a\neq1$, las funciones
\[
\log_a:(0,\infty)\to\R
\qquad\text{y}\qquad
\exp_a:\R\to(0,\infty),\quad \exp_a(x)=a^x,
\]
son inversas entre sí y satisfacen
\[
\boxed{\log_a y=x\Longleftrightarrow a^x=y}.
\]
Si $a>1$, ambas son crecientes; si $0<a<1$, ambas son decrecientes.

\subsection{Funciones trigonométricas}

Entre las funciones fundamentales se encuentran
\[
\sin x,\qquad\cos x,\qquad\tan x.
\]
Las funciones seno y coseno tienen dominio $\R$ y rango $[-1,1]$. La tangente
no está definida cuando
\[
\cos x=0,
\]
es decir,
\[
x=\frac{\pi}{2}+k\pi,\qquad k\in\mathbb Z.
\]


\subsection{Galería gráfica de las funciones elementales}

Las fórmulas anteriores deben asociarse desde el principio con una forma
gráfica. Las siguientes gráficas sirven como referencias básicas; no sustituyen
el análisis de dominio, imagen, ceros y transformaciones.

\begin{center}
\begin{tikzpicture}
\begin{axis}[width=.72\textwidth,height=5.5cm,axis lines=middle,
 xmin=-4,xmax=4,ymin=-4,ymax=6,xlabel={$x$},ylabel={$y$},grid=both,samples=150,
 legend style={font=\scriptsize,at={(0.02,0.98)},anchor=north west}]
\addplot[thick,domain=-4:4] {2}; \addlegendentry{$y=2$}
\addplot[thick,domain=-4:4] {x}; \addlegendentry{$y=x$}
\addplot[thick,domain=-3:3] {0.5*x+1}; \addlegendentry{$y=\frac12x+1$}
\addplot[thick,domain=-2.4:2.4] {x^2}; \addlegendentry{$y=x^2$}
\end{axis}
\end{tikzpicture}
\end{center}

\begin{center}
\begin{tikzpicture}
\begin{axis}[width=.72\textwidth,height=5.5cm,axis lines=middle,
 xmin=-5,xmax=5,ymin=-4,ymax=6,xlabel={$x$},ylabel={$y$},grid=both,samples=180,
 legend style={font=\scriptsize,at={(0.02,0.98)},anchor=north west}]
\addplot[thick,domain=-5:-0.15] {1/x}; \addlegendentry{$y=1/x$}
\addplot[thick,domain=0.15:5] {1/x};
\addplot[thick,domain=-5:5] {abs(x)}; \addlegendentry{$y=|x|$}
\addplot[thick,domain=0:5] {sqrt(x)}; \addlegendentry{$y=\sqrt{x}$}
\end{axis}
\end{tikzpicture}
\end{center}

\begin{center}
\begin{tikzpicture}
\begin{axis}[width=.72\textwidth,height=5.5cm,axis lines=middle,
 xmin=-4,xmax=5,ymin=-4,ymax=8,xlabel={$x$},ylabel={$y$},grid=both,samples=180,
 legend style={font=\scriptsize,at={(0.02,0.98)},anchor=north west}]
\addplot[thick,domain=-4:3] {2^x}; \addlegendentry{$y=2^x$}
\addplot[thick,domain=0.05:5] {ln(x)/ln(2)}; \addlegendentry{$y=\log_2x$}
\end{axis}
\end{tikzpicture}
\end{center}

% =========================================================
\section{Funciones definidas por partes}

\begin{definicion}
Una función definida por partes utiliza diferentes reglas de correspondencia en
diferentes subconjuntos de su dominio.
\end{definicion}

\begin{ejemplo}
Considere
\[
f(x)=
\begin{cases}
x+2,&x<0,\\
x^2,&0\leq x\leq2,\\
4,&x>2.
\end{cases}
\]
Entonces
\[
f(-3)=-1,\qquad f(1)=1,\qquad f(5)=4.
\]
En los puntos frontera debe observarse cuidadosamente qué desigualdad contiene
la igualdad.
\end{ejemplo}

% =========================================================
\section{Transformaciones de gráficas}

Sea conocida la gráfica de $y=f(x)$.

\subsection{Traslaciones verticales}

\[
g(x)=f(x)+k.
\]
Si $k>0$, la gráfica se desplaza $k$ unidades hacia arriba; si $k<0$, se
desplaza $|k|$ unidades hacia abajo.

\subsection{Traslaciones horizontales}

\[
g(x)=f(x-h).
\]
Si $h>0$, la gráfica se desplaza $h$ unidades hacia la derecha.

\begin{observacion}
El signo en las traslaciones horizontales suele producir confusión:
\[
f(x-h)\quad\text{desplaza hacia la derecha},
\qquad
f(x+h)\quad\text{hacia la izquierda}.
\]
\end{observacion}

\subsection{Reflexiones}

\[
g(x)=-f(x)
\]
refleja la gráfica respecto del eje $x$, mientras que
\[
g(x)=f(-x)
\]
la refleja respecto del eje $y$.

\subsection{Escalamientos}

Si
\[
g(x)=af(x),
\]
las ordenadas se multiplican por $a$.

Si
\[
g(x)=f(bx),
\]
el escalamiento ocurre en dirección horizontal.

\begin{ejemplo}
Partiendo de $f(x)=x^2$, la función
\[
g(x)=2(x-3)^2+1
\]
puede obtenerse mediante:
\begin{enumerate}
    \item traslado de $3$ unidades a la derecha;
    \item estiramiento vertical por factor $2$;
    \item traslado de $1$ unidad hacia arriba.
\end{enumerate}
Su vértice es $(3,1)$.
\end{ejemplo}

% =========================================================
\section{Operaciones con funciones}

Sean
\[
f:A\to\R,\qquad g:B\to\R.
\]

\begin{definicion}
Se definen:
\[
(f+g)(x)=f(x)+g(x),
\]
\[
(f-g)(x)=f(x)-g(x),
\]
\[
(fg)(x)=f(x)g(x),
\]
y, siempre que $g(x)\neq0$,
\[
\left(\frac{f}{g}\right)(x)=\frac{f(x)}{g(x)}.
\]
\end{definicion}

Para suma, diferencia y producto:
\[
\Dom(f\pm g)=\Dom(fg)=A\cap B.
\]
Para el cociente:
\[
\Dom\left(\frac fg\right)
=\{x\in A\cap B:g(x)\neq0\}.
\]

\begin{ejemplo}
Sean
\[
f(x)=\sqrt{x-1},
\qquad
g(x)=\frac{1}{x-3}.
\]
Entonces
\[
\Dom(f)=[1,\infty),
\qquad
\Dom(g)=\R\setminus\{3\}.
\]
Por tanto,
\[
\Dom(f+g)=[1,3)\cup(3,\infty).
\]
\end{ejemplo}

% =========================================================
\section{Composición de funciones}

\begin{definicion}[Composición]
Sean
\[
g:A\to B,\qquad f:B\to C.
\]
La composición de $f$ con $g$ es la función
\[
f\circ g:A\to C
\]
definida por
\[
(f\circ g)(x)=f(g(x)).
\]
\end{definicion}

Primero actúa $g$ y después actúa $f$:
\[
x\xmapsto{g}g(x)\xmapsto{f}f(g(x)).
\]

\begin{ejemplo}
Sean
\[
f(x)=x^2+1,\qquad g(x)=3x-2.
\]
Entonces
\[
(f\circ g)(x)
=f(3x-2)
=(3x-2)^2+1
=9x^2-12x+5.
\]
En cambio,
\[
(g\circ f)(x)
=g(x^2+1)
=3(x^2+1)-2
=3x^2+1.
\]
Por tanto,
\[
f\circ g\neq g\circ f.
\]
\end{ejemplo}

\subsection{Dominio de una composición}

Para que $f(g(x))$ tenga sentido deben cumplirse dos condiciones:
\[
x\in\Dom(g)
\]
y
\[
g(x)\in\Dom(f).
\]
Por tanto,
\[
\Dom(f\circ g)
=
\{x\in\Dom(g):g(x)\in\Dom(f)\}.
\]

\begin{ejemplo}
Sean
\[
f(x)=\sqrt{x},\qquad g(x)=x^2-4.
\]
Entonces
\[
(f\circ g)(x)=\sqrt{x^2-4}.
\]
Se requiere
\[
x^2-4\geq0,
\]
por lo que
\[
x\leq-2\quad\text{o}\quad x\geq2.
\]
Así,
\[
\Dom(f\circ g)=(-\infty,-2]\cup[2,\infty).
\]
\end{ejemplo}

\begin{proposicion}[Asociatividad de la composición]
Cuando las composiciones involucradas están bien definidas,
\[
(f\circ g)\circ h=f\circ(g\circ h).
\]
\end{proposicion}

\begin{proof}
Para todo $x$ del dominio correspondiente,
\[
((f\circ g)\circ h)(x)
=(f\circ g)(h(x))
=f(g(h(x))).
\]
Por otra parte,
\[
(f\circ(g\circ h))(x)
=f((g\circ h)(x))
=f(g(h(x))).
\]
Ambas funciones coinciden punto a punto.
\end{proof}

% =========================================================
\section{Funciones inyectivas, suprayectivas y biyectivas}

\begin{definicion}[Inyectividad]
Una función $f:A\to B$ es inyectiva si
\[
f(x_1)=f(x_2)\quad\Longrightarrow\quad x_1=x_2.
\]
Equivalentemente,
\[
x_1\neq x_2\quad\Longrightarrow\quad f(x_1)\neq f(x_2).
\]
\end{definicion}

\begin{definicion}[Suprayectividad]
Una función $f:A\to B$ es suprayectiva si
\[
f(A)=B.
\]
Equivalentemente,
\[
\forall y\in B\;\exists x\in A
\quad\text{tal que}\quad
f(x)=y.
\]
\end{definicion}

\begin{definicion}[Biyectividad]
Una función es biyectiva si es simultáneamente inyectiva y suprayectiva.
\end{definicion}

\begin{ejemplo}
La función
\[
f:\R\to\R,\qquad f(x)=2x-5
\]
es biyectiva.

Para probar inyectividad, supongamos
\[
f(x_1)=f(x_2).
\]
Entonces
\[
2x_1-5=2x_2-5,
\]
de donde $x_1=x_2$.

Para probar suprayectividad, sea $y\in\R$. Queremos resolver
\[
y=2x-5.
\]
Se obtiene
\[
x=\frac{y+5}{2}\in\R.
\]
Por tanto, todo $y\in\R$ tiene una preimagen.
\end{ejemplo}

\begin{ejemplo}
La función
\[
f:\R\to\R,\qquad f(x)=x^2
\]
no es inyectiva, pues
\[
f(1)=f(-1)=1.
\]
Tampoco es suprayectiva sobre $\R$, pues ningún número negativo pertenece a su
imagen.
\end{ejemplo}

% =========================================================
\section{Función inversa}

\begin{definicion}
Sea $f:A\to B$ una función biyectiva. La función inversa
\[
f^{-1}:B\to A
\]
es la función que asigna a cada $y\in B$ el único $x\in A$ tal que
\[
f(x)=y.
\]
\end{definicion}

Por definición,
\[
f^{-1}(f(x))=x,\qquad x\in A,
\]
y
\[
f(f^{-1}(y))=y,\qquad y\in B.
\]
En términos de composición,
\[
f^{-1}\circ f=\id_A,
\qquad
f\circ f^{-1}=\id_B.
\]

\begin{teorema}
Una función $f:A\to B$ admite una función inversa
\[
f^{-1}:B\to A
\]
si y sólo si $f$ es biyectiva.
\end{teorema}

\begin{proof}
Supongamos primero que existe $f^{-1}$. Si
\[
f(x_1)=f(x_2),
\]
aplicando $f^{-1}$ a ambos lados,
\[
f^{-1}(f(x_1))=f^{-1}(f(x_2)),
\]
por lo que $x_1=x_2$. Así, $f$ es inyectiva.

Además, para cada $y\in B$, tomando $x=f^{-1}(y)$ se obtiene
\[
f(x)=f(f^{-1}(y))=y.
\]
Luego $f$ es suprayectiva y, por tanto, biyectiva.

Recíprocamente, si $f$ es biyectiva, para cada $y\in B$ existe un único
$x\in A$ tal que $f(x)=y$. Definiendo
\[
f^{-1}(y)=x
\]
se obtiene una función bien definida de $B$ en $A$.
\end{proof}

\subsection{Cálculo algebraico de una inversa}

\begin{ejemplo}
Sea
\[
f(x)=3x-7.
\]
Escribimos
\[
y=3x-7.
\]
Despejamos $x$:
\[
x=\frac{y+7}{3}.
\]
Intercambiando la variable de entrada,
\[
f^{-1}(x)=\frac{x+7}{3}.
\]

Verificación:
\[
f(f^{-1}(x))
=3\left(\frac{x+7}{3}\right)-7=x.
\]
\end{ejemplo}

\begin{observacion}
La notación $f^{-1}(x)$ no significa
\[
\frac{1}{f(x)}.
\]
El recíproco de $f$ es $1/f$; la función inversa es otra noción.
\end{observacion}

% =========================================================
\section{Modelación mediante funciones}

Una de las finalidades centrales de las funciones es describir la dependencia
entre cantidades.

\subsection{Proceso general de modelación}

Un procedimiento elemental de modelación puede organizarse así:

\begin{enumerate}
    \item identificar las cantidades involucradas;
    \item elegir la variable independiente;
    \item expresar las demás cantidades en función de ella;
    \item establecer una relación matemática;
    \item determinar el dominio matemático;
    \item imponer las restricciones del problema real;
    \item interpretar el resultado en el contexto original.
\end{enumerate}

\begin{observacion}
El dominio físicamente admisible puede ser menor que el dominio algebraico de
la fórmula. Esta distinción es fundamental en problemas aplicados.
\end{observacion}

\subsection{Modelo geométrico}

\begin{ejemplo}[Rectángulo de perímetro fijo]
Un rectángulo tiene perímetro $40$ m. Exprese su área como función de la
longitud de uno de sus lados.

Sean $x$ y $y$ las longitudes de los lados. Como
\[
2x+2y=40,
\]
se obtiene
\[
y=20-x.
\]
El área es
\[
A=xy,
\]
de modo que
\[
A(x)=x(20-x)=20x-x^2.
\]

Algebraicamente la expresión está definida para todo $x\in\R$, pero las
longitudes deben ser positivas:
\[
x>0,\qquad20-x>0.
\]
Por tanto, el dominio físicamente admisible es
\[
0<x<20.
\]
\end{ejemplo}

\subsection{Modelo de costo}

\begin{ejemplo}
Una empresa tiene un costo fijo de \$5000 y un costo variable de \$120 por
unidad producida. Si $x$ representa el número de unidades, el costo total es
\[
C(x)=5000+120x.
\]
En un modelo continuo elemental puede considerarse $x\geq0$, aunque si $x$
representa unidades indivisibles, el dominio natural del problema es
\[
x\in\{0,1,2,\ldots\}.
\]
\end{ejemplo}

\subsection{Ingreso y utilidad}

Si el precio por unidad es $p(x)$ y se venden $x$ unidades, el ingreso es
\[
R(x)=xp(x).
\]
Si $C(x)$ es el costo, la utilidad se modela mediante
\[
U(x)=R(x)-C(x).
\]

\begin{ejemplo}
Suponga que el precio de venta depende de la cantidad según
\[
p(x)=500-2x
\]
y que
\[
C(x)=5000+100x.
\]
Entonces
\[
R(x)=x(500-2x)=500x-2x^2,
\]
y
\[
U(x)=R(x)-C(x)
=-2x^2+400x-5000.
\]
La función cuadrática obtenida será posteriormente susceptible de análisis
mediante derivadas para estudiar problemas de optimización.
\end{ejemplo}

\subsection{Modelo de distancia}

\begin{ejemplo}
Si un vehículo se mueve a velocidad constante de $72$ km/h, la distancia
recorrida después de $t$ horas es
\[
d(t)=72t,\qquad t\geq0.
\]
La variable independiente es el tiempo y la dependiente es la distancia.
\end{ejemplo}

\subsection{Proporcionalidad inversa}

\begin{ejemplo}
Para recorrer una distancia fija $D$, el tiempo necesario a velocidad constante
$v>0$ es
\[
T(v)=\frac{D}{v}.
\]
El tiempo es inversamente proporcional a la velocidad.
\end{ejemplo}

\subsection{Modelo exponencial elemental}

Una cantidad que cambia por un porcentaje fijo durante intervalos iguales puede
modelarse, en una primera aproximación, mediante
\[
P(t)=P_0a^t.
\]

\begin{ejemplo}
Una población inicial de $2000$ organismos aumenta $5\%$ por periodo. Entonces
\[
P(t)=2000(1.05)^t.
\]
Aquí $P_0=2000$ representa la condición inicial y $1.05$ el factor de
crecimiento por periodo.
\end{ejemplo}

% =========================================================

% =========================================================
\section{Laboratorio gráfico: funciones elementales y código en R}

Esta sección tiene un doble propósito. En primer lugar, permite reconocer
visualmente las familias de funciones que aparecerán de manera recurrente en
Cálculo Diferencial. En segundo lugar, introduce una perspectiva computacional:
la gráfica de una función puede aproximarse numéricamente evaluando la función
sobre una malla suficientemente fina de puntos de su dominio.

Para estudiantes de Ciencia de Datos y Sistemas Computacionales resulta útil
interpretar una función como una regla que transforma entradas en salidas. Si
\[
x_1,x_2,\ldots,x_n
\]
es una colección de entradas, entonces la evaluación vectorizada produce
\[
f(x_1),f(x_2),\ldots,f(x_n).
\]
En \texttt{R}, la instrucción \texttt{seq()} permite construir la malla y
\texttt{plot()} representar los pares
\[
(x_i,f(x_i)).
\]

\begin{nota}
Las gráficas generadas numéricamente no sustituyen al análisis matemático.
Una computadora dibuja una colección finita de puntos; propiedades como
dominio, rango, discontinuidades, asíntotas o monotonía deben justificarse
matemáticamente. La visualización es una herramienta de exploración y
verificación.
\end{nota}

\subsection{Plantilla general en R}

Para una función real $f$, una plantilla básica es:

\begin{verbatim}
f <- function(x) {
  # regla de correspondencia
}

x <- seq(a, b, length.out = 1000)
y <- f(x)

plot(x, y, type = "l",
     xlab = "x", ylab = "f(x)",
     main = "Grafica de f")
abline(h = 0, v = 0, lty = 3)
grid()
\end{verbatim}

El argumento \texttt{type = "l"} une mediante segmentos los puntos calculados.
Aumentar \texttt{length.out} hace más fina la discretización.

% ---------------------------------------------------------
\subsection{Función constante}

Sea
\[
f(x)=2.
\]
Su dominio es $\R$ y su imagen es $\{2\}$. La gráfica es la recta horizontal
$y=2$.

\begin{center}
\begin{tikzpicture}
\begin{axis}[
width=.72\textwidth,height=6cm,
axis lines=middle,xmin=-5,xmax=5,ymin=-1,ymax=4,
xlabel={$x$},ylabel={$f(x)$},grid=both]
\addplot[domain=-5:5,samples=2,thick] {2};
\end{axis}
\end{tikzpicture}
\end{center}

\begin{verbatim}
f <- function(x) rep(2, length(x))
x <- seq(-5, 5, length.out = 500)
plot(x, f(x), type = "l", xlab = "x", ylab = "f(x)",
     main = "Funcion constante: f(x) = 2")
abline(h = 0, v = 0, lty = 3)
grid()
\end{verbatim}

% ---------------------------------------------------------
\subsection{Función identidad}

\[
f(x)=x.
\]
Se tiene
\[
\Dom(f)=\R,\qquad \Ima(f)=\R.
\]
Es una función impar, estrictamente creciente y biyectiva de $\R$ en $\R$.

\begin{center}
\begin{tikzpicture}
\begin{axis}[width=.72\textwidth,height=6cm,axis lines=middle,
xmin=-5,xmax=5,ymin=-5,ymax=5,xlabel={$x$},ylabel={$f(x)$},grid=both]
\addplot[domain=-5:5,samples=2,thick] {x};
\end{axis}
\end{tikzpicture}
\end{center}

\begin{verbatim}
f <- function(x) x
x <- seq(-5, 5, length.out = 500)
plot(x, f(x), type = "l", asp = 1,
     xlab = "x", ylab = "f(x)",
     main = "Funcion identidad: f(x) = x")
abline(h = 0, v = 0, lty = 3)
grid()
\end{verbatim}

% ---------------------------------------------------------
\subsection{Función afín}

Consideremos
\[
f(x)=2x-1.
\]
La pendiente es $m=2$ y la ordenada al origen es $b=-1$.

\begin{center}
\begin{tikzpicture}
\begin{axis}[width=.72\textwidth,height=6cm,axis lines=middle,
xmin=-4,xmax=4,ymin=-7,ymax=7,xlabel={$x$},ylabel={$f(x)$},grid=both]
\addplot[domain=-4:4,samples=2,thick] {2*x-1};
\end{axis}
\end{tikzpicture}
\end{center}

\begin{verbatim}
f <- function(x) 2*x - 1
x <- seq(-4, 4, length.out = 500)
plot(x, f(x), type = "l",
     xlab = "x", ylab = "f(x)",
     main = "Funcion afin: f(x) = 2x - 1")
abline(h = 0, v = 0, lty = 3)
grid()
\end{verbatim}

% ---------------------------------------------------------
\subsection{Función cuadrática}

Sea
\[
f(x)=x^2-4.
\]
Su dominio es $\R$. Como $x^2\geq0$,
\[
f(x)\geq-4,
\]
de modo que
\[
\Ima(f)=[-4,\infty).
\]
Es par y su vértice es $(0,-4)$.

\begin{center}
\begin{tikzpicture}
\begin{axis}[width=.72\textwidth,height=6.5cm,axis lines=middle,
xmin=-4,xmax=4,ymin=-5,ymax=12,xlabel={$x$},ylabel={$f(x)$},grid=both]
\addplot[domain=-4:4,samples=150,thick] {x^2-4};
\end{axis}
\end{tikzpicture}
\end{center}

\begin{verbatim}
f <- function(x) x^2 - 4
x <- seq(-4, 4, length.out = 1000)
plot(x, f(x), type = "l",
     xlab = "x", ylab = "f(x)",
     main = "Funcion cuadratica: f(x) = x^2 - 4")
abline(h = 0, v = 0, lty = 3)
grid()
\end{verbatim}

% ---------------------------------------------------------
\subsection{Función cúbica}

\[
f(x)=x^3.
\]
Su dominio y su imagen son $\R$. Es impar y estrictamente creciente.

\begin{center}
\begin{tikzpicture}
\begin{axis}[width=.72\textwidth,height=6.5cm,axis lines=middle,
xmin=-3,xmax=3,ymin=-10,ymax=10,xlabel={$x$},ylabel={$f(x)$},grid=both]
\addplot[domain=-2.2:2.2,samples=150,thick] {x^3};
\end{axis}
\end{tikzpicture}
\end{center}

\begin{verbatim}
f <- function(x) x^3
x <- seq(-2.2, 2.2, length.out = 1000)
plot(x, f(x), type = "l",
     xlab = "x", ylab = "f(x)",
     main = "Funcion cubica: f(x) = x^3")
abline(h = 0, v = 0, lty = 3)
grid()
\end{verbatim}

% ---------------------------------------------------------
\subsection{Función valor absoluto}

\[
f(x)=|x|=
\begin{cases}
-x,&x<0,\\
x,&x\geq0.
\end{cases}
\]
Su dominio es $\R$, su imagen $[0,\infty)$ y es una función par.

\begin{center}
\begin{tikzpicture}
\begin{axis}[width=.72\textwidth,height=6cm,axis lines=middle,
xmin=-5,xmax=5,ymin=-1,ymax=6,xlabel={$x$},ylabel={$f(x)$},grid=both]
\addplot[domain=-5:5,samples=150,thick] {abs(x)};
\end{axis}
\end{tikzpicture}
\end{center}

\begin{verbatim}
f <- function(x) abs(x)
x <- seq(-5, 5, length.out = 1000)
plot(x, f(x), type = "l",
     xlab = "x", ylab = "f(x)",
     main = "Valor absoluto: f(x) = |x|")
abline(h = 0, v = 0, lty = 3)
grid()
\end{verbatim}

% ---------------------------------------------------------
\subsection{Función raíz cuadrada}

\[
f(x)=\sqrt{x}.
\]
La condición $x\geq0$ produce
\[
\Dom(f)=[0,\infty),\qquad \Ima(f)=[0,\infty).
\]

\begin{center}
\begin{tikzpicture}
\begin{axis}[width=.72\textwidth,height=6cm,axis lines=middle,
xmin=-1,xmax=10,ymin=-1,ymax=4,xlabel={$x$},ylabel={$f(x)$},grid=both]
\addplot[domain=0:10,samples=150,thick] {sqrt(x)};
\end{axis}
\end{tikzpicture}
\end{center}

\begin{verbatim}
f <- function(x) sqrt(x)
x <- seq(0, 10, length.out = 1000)
plot(x, f(x), type = "l",
     xlab = "x", ylab = "f(x)",
     main = "Raiz cuadrada: f(x) = sqrt(x)")
abline(h = 0, v = 0, lty = 3)
grid()
\end{verbatim}

% ---------------------------------------------------------
\subsection{Función recíproca}

\[
f(x)=\frac{1}{x}.
\]
Su dominio y rango son
\[
\R\setminus\{0\}.
\]
La función es impar. Los ejes coordenados son asíntotas:
\[
x=0,\qquad y=0.
\]

No debe generarse una única secuencia que atraviese $x=0$ y unirla con líneas,
pues ello puede producir una representación visual engañosa. Deben graficarse
las dos ramas por separado.

\begin{center}
\begin{tikzpicture}
\begin{axis}[width=.72\textwidth,height=6.5cm,axis lines=middle,
xmin=-5,xmax=5,ymin=-5,ymax=5,xlabel={$x$},ylabel={$f(x)$},grid=both,
restrict y to domain=-5:5]
\addplot[domain=-5:-0.2,samples=150,thick] {1/x};
\addplot[domain=0.2:5,samples=150,thick] {1/x};
\end{axis}
\end{tikzpicture}
\end{center}

\begin{verbatim}
f <- function(x) 1/x

x1 <- seq(-5, -0.05, length.out = 1000)
x2 <- seq(0.05, 5, length.out = 1000)

plot(x1, f(x1), type = "l",
     xlim = c(-5, 5), ylim = c(-5, 5),
     xlab = "x", ylab = "f(x)",
     main = "Funcion reciproca: f(x) = 1/x")
lines(x2, f(x2))
abline(h = 0, v = 0, lty = 3)
grid()
\end{verbatim}

% ---------------------------------------------------------
\subsection{Función racional con traslación}

Sea
\[
f(x)=\frac{x+1}{x-2}.
\]
Puede reescribirse como
\[
f(x)=1+\frac{3}{x-2}.
\]
Esta forma permite reconocer las asíntotas
\[
x=2,\qquad y=1.
\]

\begin{center}
\begin{tikzpicture}
\begin{axis}[width=.72\textwidth,height=6.5cm,axis lines=middle,
xmin=-5,xmax=7,ymin=-6,ymax=7,xlabel={$x$},ylabel={$f(x)$},grid=both,
restrict y to domain=-6:7]
\addplot[domain=-5:1.9,samples=180,thick] {(x+1)/(x-2)};
\addplot[domain=2.1:7,samples=180,thick] {(x+1)/(x-2)};
\addplot[dashed] coordinates {(2,-6) (2,7)};
\addplot[dashed] coordinates {(-5,1) (7,1)};
\end{axis}
\end{tikzpicture}
\end{center}

\begin{verbatim}
f <- function(x) (x + 1)/(x - 2)

x1 <- seq(-5, 1.95, length.out = 1000)
x2 <- seq(2.05, 7, length.out = 1000)

plot(x1, f(x1), type = "l",
     xlim = c(-5, 7), ylim = c(-6, 7),
     xlab = "x", ylab = "f(x)",
     main = "Funcion racional")
lines(x2, f(x2))
abline(v = 2, h = 1, lty = 2)
abline(h = 0, v = 0, lty = 3)
grid()
\end{verbatim}

% ---------------------------------------------------------
\subsection{Función exponencial}

Sea
\[
f(x)=2^x.
\]
Entonces
\[
\Dom(f)=\R,\qquad \Ima(f)=(0,\infty).
\]
Además,
\[
f(0)=1.
\]
La función es estrictamente creciente.

\begin{center}
\begin{tikzpicture}
\begin{axis}[width=.72\textwidth,height=6cm,axis lines=middle,
xmin=-4,xmax=4,ymin=-1,ymax=10,xlabel={$x$},ylabel={$f(x)$},grid=both]
\addplot[domain=-4:3.2,samples=180,thick] {2^x};
\end{axis}
\end{tikzpicture}
\end{center}

\begin{verbatim}
f <- function(x) 2^x
x <- seq(-4, 3.2, length.out = 1000)
plot(x, f(x), type = "l",
     xlab = "x", ylab = "f(x)",
     main = "Funcion exponencial: f(x) = 2^x")
abline(h = 0, v = 0, lty = 3)
grid()
\end{verbatim}

% ---------------------------------------------------------
\subsection{Función logaritmo natural}

\[
f(x)=\ln x.
\]
Su dominio e imagen son
\[
\Dom(f)=(0,\infty),
\qquad
\Ima(f)=\R.
\]
Además,
\[
\ln(1)=0.
\]

\begin{center}
\begin{tikzpicture}
\begin{axis}[width=.72\textwidth,height=6cm,axis lines=middle,
xmin=-1,xmax=10,ymin=-4,ymax=4,xlabel={$x$},ylabel={$f(x)$},grid=both]
\addplot[domain=0.03:10,samples=200,thick] {ln(x)};
\end{axis}
\end{tikzpicture}
\end{center}

\begin{verbatim}
f <- function(x) log(x)
x <- seq(0.03, 10, length.out = 1000)
plot(x, f(x), type = "l",
     xlab = "x", ylab = "f(x)",
     main = "Logaritmo natural: f(x) = ln(x)")
abline(h = 0, v = 0, lty = 3)
grid()
\end{verbatim}

\begin{nota}
En \texttt{R}, \texttt{log(x)} representa el logaritmo natural. Para una base
$b$ puede utilizarse \texttt{log(x, base = b)}.
\end{nota}

% ---------------------------------------------------------
\subsection{Funciones seno y coseno}

\[
f(x)=\sin x,\qquad g(x)=\cos x.
\]
Ambas tienen dominio $\R$, imagen $[-1,1]$ y periodo $2\pi$. El seno es impar y
el coseno es par.

\begin{center}
\begin{tikzpicture}
\begin{axis}[width=.82\textwidth,height=6cm,axis lines=middle,
xmin=-6.4,xmax=6.4,ymin=-1.5,ymax=1.5,xlabel={$x$},ylabel={$y$},grid=both,
legend style={at={(0.5,-0.18)},anchor=north,legend columns=2}]
\addplot[domain=-2*pi:2*pi,samples=250,thick] {sin(deg(x))};
\addlegendentry{$\sin x$}
\addplot[domain=-2*pi:2*pi,samples=250,thick,dashed] {cos(deg(x))};
\addlegendentry{$\cos x$}
\end{axis}
\end{tikzpicture}
\end{center}

\begin{verbatim}
x <- seq(-2*pi, 2*pi, length.out = 1500)

plot(x, sin(x), type = "l",
     ylim = c(-1.5, 1.5),
     xlab = "x", ylab = "y",
     main = "Funciones seno y coseno")
lines(x, cos(x), lty = 2)
abline(h = 0, v = 0, lty = 3)
legend("topright",
       legend = c("sin(x)", "cos(x)"),
       lty = c(1, 2))
grid()
\end{verbatim}

% ---------------------------------------------------------
\subsection{Función tangente}

\[
f(x)=\tan x=\frac{\sin x}{\cos x}.
\]
No está definida cuando
\[
x=\frac{\pi}{2}+k\pi,\qquad k\in\mathbb Z.
\]
Por ello, nuevamente deben evitarse uniones artificiales a través de las
asíntotas verticales.

\begin{center}
\begin{tikzpicture}
\begin{axis}[width=.82\textwidth,height=6.5cm,axis lines=middle,
xmin=-4.7,xmax=4.7,ymin=-5,ymax=5,xlabel={$x$},ylabel={$f(x)$},grid=both,
restrict y to domain=-5:5]
\addplot[domain=-4.6:-1.62,samples=150,thick] {tan(deg(x))};
\addplot[domain=-1.52:1.52,samples=150,thick] {tan(deg(x))};
\addplot[domain=1.62:4.6,samples=150,thick] {tan(deg(x))};
\end{axis}
\end{tikzpicture}
\end{center}

\begin{verbatim}
f <- function(x) tan(x)

# Se grafican intervalos separados por las asintotas.
eps <- 0.03
intervalos <- list(
  c(-3*pi/2 + eps, -pi/2 - eps),
  c(-pi/2 + eps,   pi/2 - eps),
  c( pi/2 + eps, 3*pi/2 - eps)
)

plot(NA, xlim = c(-3*pi/2, 3*pi/2), ylim = c(-5, 5),
     xlab = "x", ylab = "tan(x)",
     main = "Funcion tangente")

for (I in intervalos) {
  x <- seq(I[1], I[2], length.out = 500)
  lines(x, f(x))
}

abline(v = c(-3*pi/2, -pi/2, pi/2, 3*pi/2), lty = 2)
abline(h = 0, v = 0, lty = 3)
grid()
\end{verbatim}

% ---------------------------------------------------------
\subsection{Función definida por partes}

Considere
\[
f(x)=
\begin{cases}
-x-1,&x<0,\\
x^2,&0\leq x\leq2,\\
4,&x>2.
\end{cases}
\]

La definición por partes se implementa en R mediante \texttt{ifelse}. Para más
de dos casos pueden anidarse condiciones.

\begin{center}
\begin{tikzpicture}
\begin{axis}[width=.75\textwidth,height=6.5cm,axis lines=middle,
xmin=-4,xmax=5,ymin=-2,ymax=6,xlabel={$x$},ylabel={$f(x)$},grid=both]
\addplot[domain=-4:-0.02,samples=80,thick] {-x-1};
\addplot[domain=0:2,samples=80,thick] {x^2};
\addplot[domain=2.02:5,samples=2,thick] {4};
\addplot[only marks,mark=o] coordinates {(0,-1)};
\addplot[only marks,mark=*] coordinates {(0,0) (2,4)};
\end{axis}
\end{tikzpicture}
\end{center}

\begin{verbatim}
f <- function(x) {
  ifelse(x < 0,
         -x - 1,
         ifelse(x <= 2, x^2, 4))
}

x <- seq(-4, 5, length.out = 1500)
plot(x, f(x), type = "l",
     xlab = "x", ylab = "f(x)",
     main = "Funcion definida por partes")
points(0, -1, pch = 1)  # punto abierto
points(0,  0, pch = 19) # punto cerrado
points(2,  4, pch = 19)
abline(h = 0, v = 0, lty = 3)
grid()
\end{verbatim}

% ---------------------------------------------------------
\subsection{Transformaciones de una función madre}

Partamos de
\[
f(x)=x^2
\]
y consideremos
\[
g(x)=2(x-3)^2+1.
\]
La comparación gráfica permite visualizar simultáneamente la traslación
horizontal, el escalamiento vertical y la traslación vertical.

\begin{center}
\begin{tikzpicture}
\begin{axis}[width=.78\textwidth,height=6.5cm,axis lines=middle,
xmin=-4,xmax=7,ymin=-1,ymax=15,xlabel={$x$},ylabel={$y$},grid=both,
legend style={at={(0.5,-0.18)},anchor=north,legend columns=2}]
\addplot[domain=-3.5:3.5,samples=150,thick,dashed] {x^2};
\addlegendentry{$f(x)=x^2$}
\addplot[domain=0.4:5.6,samples=150,thick] {2*(x-3)^2+1};
\addlegendentry{$g(x)=2(x-3)^2+1$}
\end{axis}
\end{tikzpicture}
\end{center}

\begin{verbatim}
f <- function(x) x^2
g <- function(x) 2*(x - 3)^2 + 1

x <- seq(-4, 7, length.out = 1500)

plot(x, f(x), type = "l",
     ylim = c(-1, 15),
     xlab = "x", ylab = "y",
     main = "Transformacion de f(x) = x^2")
lines(x, g(x), lty = 2)
legend("topright",
       legend = c("f(x) = x^2", "g(x) = 2(x-3)^2 + 1"),
       lty = c(1, 2))
abline(h = 0, v = 0, lty = 3)
grid()
\end{verbatim}

% ---------------------------------------------------------
\subsection{Composición: visualización computacional}

Sean
\[
f(x)=x^2+1,\qquad g(x)=2x-3.
\]
Entonces
\[
(f\circ g)(x)=(2x-3)^2+1
\]
y
\[
(g\circ f)(x)=2(x^2+1)-3=2x^2-1.
\]
La comparación gráfica refuerza que la composición no es conmutativa.

\begin{center}
\begin{tikzpicture}
\begin{axis}[width=.78\textwidth,height=6.5cm,axis lines=middle,
xmin=-2,xmax=4,ymin=-3,ymax=18,xlabel={$x$},ylabel={$y$},grid=both,
legend style={at={(0.5,-0.18)},anchor=north,legend columns=2}]
\addplot[domain=-1.5:4,samples=180,thick] {(2*x-3)^2+1};
\addlegendentry{$f\circ g$}
\addplot[domain=-1.5:3,samples=180,thick,dashed] {2*x^2-1};
\addlegendentry{$g\circ f$}
\end{axis}
\end{tikzpicture}
\end{center}

\begin{verbatim}
f <- function(x) x^2 + 1
g <- function(x) 2*x - 3

fog <- function(x) f(g(x))
gof <- function(x) g(f(x))

x <- seq(-1.5, 4, length.out = 1200)

plot(x, fog(x), type = "l",
     ylim = c(-3, 18),
     xlab = "x", ylab = "y",
     main = "La composicion no es conmutativa")
lines(x, gof(x), lty = 2)
legend("topright",
       legend = c("(f o g)(x)", "(g o f)(x)"),
       lty = c(1, 2))
abline(h = 0, v = 0, lty = 3)
grid()
\end{verbatim}

% ---------------------------------------------------------
\subsection{Función e inversa: simetría respecto de $y=x$}

Sea
\[
f(x)=2x-3.
\]
Su inversa es
\[
f^{-1}(x)=\frac{x+3}{2}.
\]
Las gráficas de $f$ y $f^{-1}$ son reflejos una de la otra respecto de la
recta
\[
y=x.
\]

\begin{center}
\begin{tikzpicture}
\begin{axis}[width=.72\textwidth,height=7cm,axis lines=middle,
xmin=-5,xmax=5,ymin=-5,ymax=5,xlabel={$x$},ylabel={$y$},grid=both,
legend style={at={(0.5,-0.18)},anchor=north,legend columns=3}]
\addplot[domain=-4:4,samples=2,thick] {2*x-3};
\addlegendentry{$f$}
\addplot[domain=-5:5,samples=2,thick,dashed] {(x+3)/2};
\addlegendentry{$f^{-1}$}
\addplot[domain=-5:5,samples=2,dotted] {x};
\addlegendentry{$y=x$}
\end{axis}
\end{tikzpicture}
\end{center}

\begin{verbatim}
f    <- function(x) 2*x - 3
finv <- function(x) (x + 3)/2

x <- seq(-5, 5, length.out = 500)

plot(x, f(x), type = "l",
     xlim = c(-5, 5), ylim = c(-5, 5), asp = 1,
     xlab = "x", ylab = "y",
     main = "Funcion e inversa")
lines(x, finv(x), lty = 2)
abline(a = 0, b = 1, lty = 3)
legend("topleft",
       legend = c("f", "f^{-1}", "y = x"),
       lty = c(1, 2, 3))
grid()
\end{verbatim}

% ---------------------------------------------------------
\subsection{Ejercicio computacional guiado}

Considere
\[
f(x)=\frac{\sqrt{x+2}}{x-1}.
\]

\begin{enumerate}
    \item Determine analíticamente el dominio.
    \item Identifique el punto excluido debido al denominador.
    \item Construya en R dos mallas, una a cada lado de $x=1$.
    \item Grafique ambas ramas sin unirlas.
    \item Compare la gráfica obtenida con el dominio calculado.
    \item Experimente con valores cada vez más próximos a $x=1$.
    \item Explique numéricamente qué observa, sin utilizar todavía la definición
    formal de límite.
\end{enumerate}

Código inicial:

\begin{verbatim}
f <- function(x) sqrt(x + 2)/(x - 1)

x1 <- seq(-2, 0.99, length.out = 1000)
x2 <- seq(1.01, 8, length.out = 1000)

plot(x1, f(x1), type = "l",
     xlim = c(-2, 8), ylim = c(-20, 20),
     xlab = "x", ylab = "f(x)",
     main = "Exploracion de una funcion racional-radical")
lines(x2, f(x2))
abline(v = 1, lty = 2)
abline(h = 0, v = 0, lty = 3)
grid()
\end{verbatim}

\subsection{Actividad para Ciencia de Datos y Sistemas Computacionales}

Construya un pequeño catálogo computacional de funciones. Para cada función,
genere una tabla con $n=1001$ puntos y las columnas
\[
x,\qquad f(x),
\]
y obtenga su gráfica.

En R:

\begin{verbatim}
f <- function(x) x^2 - 4
x <- seq(-5, 5, length.out = 1001)

datos <- data.frame(
  x = x,
  fx = f(x)
)

head(datos)
summary(datos)

plot(datos$x, datos$fx, type = "l",
     xlab = "x", ylab = "f(x)")
grid()
\end{verbatim}

Discuta las siguientes preguntas:

\begin{enumerate}
    \item ¿Una tabla finita permite conocer exactamente la imagen de una
    función definida sobre un intervalo?
    \item ¿Qué información matemática se pierde al discretizar el dominio?
    \item ¿Qué ocurre cerca de una discontinuidad si la malla es demasiado
    gruesa?
    \item ¿Cómo influye \texttt{length.out} en la apariencia de la gráfica?
    \item ¿Puede una gráfica numérica constituir una demostración?
    \item ¿Qué ventajas aporta, entonces, el análisis computacional?
\end{enumerate}

La respuesta a la quinta pregunta es particularmente importante: una
visualización puede sugerir una propiedad y aportar evidencia numérica, pero
no sustituye una demostración cuando se requiere una afirmación válida para
todos los elementos de un conjunto infinito.


% =========================================================
\section{Banco de ejemplos para elegir en clase}
\label{sec:banco-ejemplos}
% =========================================================

Esta sección está pensada como un \emph{banco de trabajo para el pizarrón}. No es necesario resolver todos los ejemplos en una misma sesión. El profesor puede seleccionar, según el avance del grupo, un ejemplo básico, uno intermedio o uno de reto. En todos los casos conviene pedir primero al estudiante que identifique qué se pregunta, qué restricciones aparecen y qué información de la función puede obtenerse antes de efectuar cálculos.

\subsection{Funciones lineales y afines}

\begin{ejemplo}[Básico: evaluación e interpretación]
Sea $f(x)=3x-5$. Calcule $f(0)$, $f(2)$ y determine la intersección con el eje $x$.

\textbf{Solución.} Sustituyendo,
\[
f(0)=-5,\qquad f(2)=1.
\]
Para hallar la intersección con el eje $x$ imponemos $f(x)=0$:
\[
3x-5=0\quad\Longrightarrow\quad x=\frac53.
\]
Por tanto, la gráfica corta a los ejes en $(0,-5)$ y $(5/3,0)$. La pendiente es $3>0$, así que la función es creciente.


\textbf{Gráfica de apoyo.}\medskip
\begin{center}
\begin{tikzpicture}
\begin{axis}[
    width=0.78\textwidth,height=5.4cm,
    axis lines=middle,xmin=-1,xmax=3,ymin=-6,ymax=5,
    xlabel={$x$},ylabel={$y$},grid=both,
    minor tick num=1,samples=180,
    legend style={font=\scriptsize,at={(0.02,0.98)},anchor=north west}
]
\addplot[thick,domain=-1:3] {3*x-5};
\addplot[only marks,mark=*] coordinates {(0,-5) (2,1) (1.6667,0)};
\end{axis}
\end{tikzpicture}
\end{center}

\end{ejemplo}

\begin{ejemplo}[Intermedio: recta determinada por dos puntos]
Determine la función afín cuya gráfica pasa por $P=(-2,5)$ y $Q=(4,-1)$.

\textbf{Solución.} La pendiente es
\[
m=\frac{-1-5}{4-(-2)}=\frac{-6}{6}=-1.
\]
Usando la forma punto--pendiente con $P$,
\[
y-5=-(x+2),
\]
de donde
\[
\boxed{f(x)=-x+3}.
\]
La comprobación es inmediata: $f(-2)=5$ y $f(4)=-1$.


\textbf{Gráfica de apoyo.}\medskip
\begin{center}
\begin{tikzpicture}
\begin{axis}[
    width=0.78\textwidth,height=5.4cm,
    axis lines=middle,xmin=-4,xmax=6,ymin=-4,ymax=8,
    xlabel={$x$},ylabel={$y$},grid=both,
    minor tick num=1,samples=180,
    legend style={font=\scriptsize,at={(0.02,0.98)},anchor=north west}
]
\addplot[thick,domain=-3:5] {-x+3};
\addplot[only marks,mark=*] coordinates {(-2,5) (4,-1)};
\end{axis}
\end{tikzpicture}
\end{center}

\end{ejemplo}

\begin{ejemplo}[Intermedio: parámetro]
Sea $f(x)=(k-2)x+4$. Determine $k$ para que la gráfica sea paralela a $g(x)=5x-1$ y determine $k$ para que $f$ sea constante.

\textbf{Solución.} Rectas paralelas tienen la misma pendiente, por lo que
\[
k-2=5\quad\Longrightarrow\quad \boxed{k=7}.
\]
Para que $f$ sea constante su pendiente debe ser cero:
\[
k-2=0\quad\Longrightarrow\quad \boxed{k=2}.
\]


\textbf{Gráfica de apoyo.}\medskip
\begin{center}
\begin{tikzpicture}
\begin{axis}[
    width=0.78\textwidth,height=5.4cm,
    axis lines=middle,xmin=-2,xmax=2,ymin=-8,ymax=10,
    xlabel={$x$},ylabel={$y$},grid=both,
    minor tick num=1,samples=180,
    legend style={font=\scriptsize,at={(0.02,0.98)},anchor=north west}
]
\addplot[thick,domain=-2:2] {5*x+4};
\addplot[thick,dashed,domain=-2:2] {5*x-1};
\addplot[thick,dotted,domain=-2:2] {4};
\end{axis}
\end{tikzpicture}
\end{center}

\end{ejemplo}

\begin{ejemplo}[Reto: construcción de un modelo lineal]
Una empresa cobra una cuota fija de $120$ pesos y $18$ pesos por cada unidad consumida. Construya la función de costo y determine el consumo correspondiente a un cobro de $480$ pesos.

\textbf{Solución.} Si $x$ es el número de unidades consumidas,
\[
C(x)=120+18x,\qquad x\ge0.
\]
Ahora imponemos $C(x)=480$:
\[
120+18x=480\Longrightarrow 18x=360\Longrightarrow \boxed{x=20}.
\]
El término $120$ representa el costo fijo y $18$ la razón de cambio del costo respecto del consumo.


\textbf{Gráfica de apoyo.}\medskip
\begin{center}
\begin{tikzpicture}
\begin{axis}[
    width=0.78\textwidth,height=5.4cm,
    axis lines=middle,xmin=0,xmax=25,ymin=0,ymax=600,
    xlabel={$x$},ylabel={$y$},grid=both,
    minor tick num=1,samples=180,
    legend style={font=\scriptsize,at={(0.02,0.98)},anchor=north west}
]
\addplot[thick,domain=0:25] {120+18*x};
\addplot[only marks,mark=*] coordinates {(0,120) (20,480)};
\end{axis}
\end{tikzpicture}
\end{center}

\end{ejemplo}

\subsection{Funciones cuadráticas y polinomiales}

\begin{ejemplo}[Básico: ceros y factorización]
Analice los ceros de $f(x)=x^2-5x+6$.

\textbf{Solución.} Factorizamos:
\[
f(x)=(x-2)(x-3).
\]
Así,
\[
f(x)=0\Longleftrightarrow x=2\ \text{o}\ x=3.
\]
La gráfica corta al eje $x$ en $(2,0)$ y $(3,0)$. Como el coeficiente de $x^2$ es positivo, la parábola abre hacia arriba.


\textbf{Gráfica de apoyo.}\medskip
\begin{center}
\begin{tikzpicture}
\begin{axis}[
    width=0.78\textwidth,height=5.4cm,
    axis lines=middle,xmin=-1,xmax=7,ymin=-3,ymax=12,
    xlabel={$x$},ylabel={$y$},grid=both,
    minor tick num=1,samples=180,
    legend style={font=\scriptsize,at={(0.02,0.98)},anchor=north west}
]
\addplot[thick,domain=-1:7] {x^2-5*x+6};
\addplot[only marks,mark=*] coordinates {(2,0) (3,0)};
\end{axis}
\end{tikzpicture}
\end{center}

\end{ejemplo}

\begin{ejemplo}[Intermedio: completar el cuadrado]
Determine el vértice y la imagen de $f(x)=x^2-6x+5$.

\textbf{Solución.}
\[
x^2-6x+5=(x-3)^2-9+5=(x-3)^2-4.
\]
Por tanto,
\[
\boxed{f(x)=(x-3)^2-4}.
\]
El vértice es $(3,-4)$. Como $(x-3)^2\ge0$,
\[
f(x)\ge-4,
\]
y por ello
\[
\boxed{\Ima(f)=[-4,\infty)}.
\]


\textbf{Gráfica de apoyo.}\medskip
\begin{center}
\begin{tikzpicture}
\begin{axis}[
    width=0.78\textwidth,height=5.4cm,
    axis lines=middle,xmin=-1,xmax=7,ymin=-6,ymax=12,
    xlabel={$x$},ylabel={$y$},grid=both,
    minor tick num=1,samples=180,
    legend style={font=\scriptsize,at={(0.02,0.98)},anchor=north west}
]
\addplot[thick,domain=-1:7] {(x-3)^2-4};
\addplot[only marks,mark=*] coordinates {(3,-4)};
\end{axis}
\end{tikzpicture}
\end{center}

\end{ejemplo}

\begin{ejemplo}[Intermedio: simetría de un polinomio]
Determine si $f(x)=x^4-3x^2+2$ es par, impar o ninguna.

\textbf{Solución.}
\[
f(-x)=(-x)^4-3(-x)^2+2=x^4-3x^2+2=f(x).
\]
Luego $f$ es \textbf{par}; su gráfica es simétrica respecto del eje $y$.


\textbf{Gráfica de apoyo.}\medskip
\begin{center}
\begin{tikzpicture}
\begin{axis}[
    width=0.78\textwidth,height=5.4cm,
    axis lines=middle,xmin=-2.5,xmax=2.5,ymin=-2,ymax=8,
    xlabel={$x$},ylabel={$y$},grid=both,
    minor tick num=1,samples=180,
    legend style={font=\scriptsize,at={(0.02,0.98)},anchor=north west}
]
\addplot[thick,domain=-2.2:2.2] {x^4-3*x^2+2};
\end{axis}
\end{tikzpicture}
\end{center}

\end{ejemplo}

\begin{ejemplo}[Reto: lectura estructural]
Para $f(x)=x^3-4x=x(x-2)(x+2)$ determine dominio, ceros, paridad y signo en los intervalos determinados por sus ceros.

\textbf{Solución.} Al ser polinomial, $\Dom(f)=\R$. Sus ceros son $-2,0,2$. Además,
\[
f(-x)=-x^3+4x=-f(x),
\]
por lo que es impar. El análisis de signos de los factores produce
\[
f(x)<0\text{ en }(-\infty,-2)\cup(0,2),
\qquad
f(x)>0\text{ en }(-2,0)\cup(2,\infty).
\]


\textbf{Gráfica de apoyo.}\medskip
\begin{center}
\begin{tikzpicture}
\begin{axis}[
    width=0.78\textwidth,height=5.4cm,
    axis lines=middle,xmin=-3,xmax=3,ymin=-8,ymax=8,
    xlabel={$x$},ylabel={$y$},grid=both,
    minor tick num=1,samples=180,
    legend style={font=\scriptsize,at={(0.02,0.98)},anchor=north west}
]
\addplot[thick,domain=-3:3] {x^3-4*x};
\addplot[only marks,mark=*] coordinates {(-2,0) (0,0) (2,0)};
\end{axis}
\end{tikzpicture}
\end{center}

\end{ejemplo}

\subsection{Funciones racionales}

\begin{ejemplo}[Básico: dominio]
Determine el dominio de
\[
f(x)=\frac{x+1}{x-4}.
\]
\textbf{Solución.} El denominador no puede anularse:
\[
x-4\neq0\Longrightarrow x\neq4.
\]
Por tanto,
\[
\boxed{\Dom(f)=\R\setminus\{4\}}.
\]


\textbf{Gráfica de apoyo.}\medskip
\begin{center}
\begin{tikzpicture}
\begin{axis}[
    width=0.78\textwidth,height=5.4cm,
    axis lines=middle,xmin=-5,xmax=8,ymin=-6,ymax=8,
    xlabel={$x$},ylabel={$y$},grid=both,
    minor tick num=1,samples=180,
    legend style={font=\scriptsize,at={(0.02,0.98)},anchor=north west}
]
\addplot[thick,domain=-5:3.9] {(x+1)/(x-4)};
\addplot[thick,domain=4.1:8] {(x+1)/(x-4)};
\addplot[dashed] coordinates {(4,-6) (4,8)};
\end{axis}
\end{tikzpicture}
\end{center}

\end{ejemplo}

\begin{ejemplo}[Intermedio: simplificar no recupera el dominio]
Analice
\[
f(x)=\frac{x^2-4}{x-2}.
\]
\textbf{Solución.} Antes de simplificar debemos registrar la restricción $x\neq2$. Después,
\[
f(x)=\frac{(x-2)(x+2)}{x-2}=x+2,\qquad x\neq2.
\]
La gráfica coincide con la recta $y=x+2$ salvo por el punto $(2,4)$, que no pertenece a la gráfica. Así,
\[
\boxed{\Dom(f)=\R\setminus\{2\}}.
\]
Este ejemplo muestra que una simplificación algebraica no modifica las restricciones heredadas de la expresión original.


\textbf{Gráfica de apoyo.}\medskip
\begin{center}
\begin{tikzpicture}
\begin{axis}[
    width=0.78\textwidth,height=5.4cm,
    axis lines=middle,xmin=-4,xmax=6,ymin=-3,ymax=9,
    xlabel={$x$},ylabel={$y$},grid=both,
    minor tick num=1,samples=180,
    legend style={font=\scriptsize,at={(0.02,0.98)},anchor=north west}
]
\addplot[thick,domain=-4:1.98] {x+2};
\addplot[thick,domain=2.02:6] {x+2};
\addplot[only marks,mark=o,mark options={fill=white},very thick] coordinates {(2,4)};
\end{axis}
\end{tikzpicture}
\end{center}

\end{ejemplo}

\begin{ejemplo}[Intermedio: varias restricciones]
Determine el dominio de
\[
f(x)=\frac{2x-1}{x^2+x-6}.
\]
\textbf{Solución.} Factorizamos el denominador:
\[
x^2+x-6=(x+3)(x-2).
\]
Por tanto,
\[
x\neq-3,\qquad x\neq2,
\]
y
\[
\boxed{\Dom(f)=\R\setminus\{-3,2\}}.
\]


\textbf{Gráfica de apoyo.}\medskip
\begin{center}
\begin{tikzpicture}
\begin{axis}[
    width=0.78\textwidth,height=5.4cm,
    axis lines=middle,xmin=-7,xmax=7,ymin=-6,ymax=6,
    xlabel={$x$},ylabel={$y$},grid=both,
    minor tick num=1,samples=180,
    legend style={font=\scriptsize,at={(0.02,0.98)},anchor=north west}
]
\addplot[thick,domain=-7:-3.08] {(2*x-1)/(x^2+x-6)};
\addplot[thick,domain=-2.92:1.92] {(2*x-1)/(x^2+x-6)};
\addplot[thick,domain=2.08:7] {(2*x-1)/(x^2+x-6)};
\addplot[dashed] coordinates {(-3,-6) (-3,6)};
\addplot[dashed] coordinates {(2,-6) (2,6)};
\end{axis}
\end{tikzpicture}
\end{center}

\end{ejemplo}

\begin{ejemplo}[Intermedio: imagen mediante despeje]
Determine la imagen de
\[
f(x)=\frac{2x+1}{x-3}.
\]
\textbf{Solución.} Escribimos
\[
y=\frac{2x+1}{x-3}
\]
y despejamos $x$:
\[
yx-3y=2x+1,
\qquad x(y-2)=3y+1,
\qquad x=\frac{3y+1}{y-2}.
\]
Este despeje es posible para todo $y$ excepto $y=2$. Luego
\[
\boxed{\Ima(f)=\R\setminus\{2\}}.
\]


\textbf{Gráfica de apoyo.}\medskip
\begin{center}
\begin{tikzpicture}
\begin{axis}[
    width=0.78\textwidth,height=5.4cm,
    axis lines=middle,xmin=-8,xmax=10,ymin=-8,ymax=10,
    xlabel={$x$},ylabel={$y$},grid=both,
    minor tick num=1,samples=180,
    legend style={font=\scriptsize,at={(0.02,0.98)},anchor=north west}
]
\addplot[thick,domain=-8:2.92] {(2*x+1)/(x-3)};
\addplot[thick,domain=3.08:10] {(2*x+1)/(x-3)};
\addplot[dashed] coordinates {(3,-8) (3,10)};
\addplot[dashed,domain=-8:10] {2};
\end{axis}
\end{tikzpicture}
\end{center}

\end{ejemplo}

\begin{ejemplo}[Reto: dominio de una racional compuesta]
Determine el dominio de
\[
f(x)=\frac{x+2}{(x-1)(x^2-9)}.
\]
\textbf{Solución.} Los valores prohibidos son las raíces del denominador:
\[
x-1=0,\qquad x^2-9=(x-3)(x+3)=0.
\]
Por tanto,
\[
\boxed{\Dom(f)=\R\setminus\{-3,1,3\}}.
\]


\textbf{Gráfica de apoyo.}\medskip
\begin{center}
\begin{tikzpicture}
\begin{axis}[
    width=0.78\textwidth,height=5.4cm,
    axis lines=middle,xmin=-6,xmax=6,ymin=-4,ymax=4,
    xlabel={$x$},ylabel={$y$},grid=both,
    minor tick num=1,samples=180,
    legend style={font=\scriptsize,at={(0.02,0.98)},anchor=north west}
]
\addplot[thick,domain=-6:-3.08] {(x+2)/((x-1)*(x^2-9))};
\addplot[thick,domain=-2.92:0.92] {(x+2)/((x-1)*(x^2-9))};
\addplot[thick,domain=1.08:2.92] {(x+2)/((x-1)*(x^2-9))};
\addplot[thick,domain=3.08:6] {(x+2)/((x-1)*(x^2-9))};
\foreach \a in {-3,1,3}{\addplot[dashed] coordinates {(\a,-4) (\a,4)};}
\end{axis}
\end{tikzpicture}
\end{center}

\end{ejemplo}

\subsection{Funciones radicales}

\begin{ejemplo}[Básico: raíz cuadrada]
Determine el dominio de $f(x)=\sqrt{x-3}$.

\textbf{Solución.} Para una raíz de índice par exigimos
\[
x-3\ge0,
\]
de modo que
\[
\boxed{\Dom(f)=[3,\infty)}.
\]


\textbf{Gráfica de apoyo.}\medskip
\begin{center}
\begin{tikzpicture}
\begin{axis}[
    width=0.78\textwidth,height=5.4cm,
    axis lines=middle,xmin=0,xmax=9,ymin=-1,ymax=4,
    xlabel={$x$},ylabel={$y$},grid=both,
    minor tick num=1,samples=180,
    legend style={font=\scriptsize,at={(0.02,0.98)},anchor=north west}
]
\addplot[thick,domain=3:9] {sqrt(x-3)};
\addplot[only marks,mark=*] coordinates {(3,0)};
\end{axis}
\end{tikzpicture}
\end{center}

\end{ejemplo}

\begin{ejemplo}[Intermedio: radical cuadrático]
Determine el dominio y la imagen de
\[
f(x)=\sqrt{9-x^2}.
\]
\textbf{Solución.} Exigimos
\[
9-x^2\ge0\Longleftrightarrow x^2\le9\Longleftrightarrow -3\le x\le3.
\]
Así,
\[
\boxed{\Dom(f)=[-3,3]}.
\]
Además $0\le\sqrt{9-x^2}\le3$, por lo que
\[
\boxed{\Ima(f)=[0,3]}.
\]
Geométricamente es la semicircunferencia superior $x^2+y^2=9$.


\textbf{Gráfica de apoyo.}\medskip
\begin{center}
\begin{tikzpicture}
\begin{axis}[
    width=0.78\textwidth,height=5.4cm,
    axis lines=middle,xmin=-4,xmax=4,ymin=-1,ymax=4,
    xlabel={$x$},ylabel={$y$},grid=both,
    minor tick num=1,samples=180,
    legend style={font=\scriptsize,at={(0.02,0.98)},anchor=north west}
]
\addplot[thick,domain=-3:3] {sqrt(9-x^2)};
\addplot[only marks,mark=*] coordinates {(-3,0) (0,3) (3,0)};
\end{axis}
\end{tikzpicture}
\end{center}

\end{ejemplo}

\begin{ejemplo}[Intermedio: cociente dentro de una raíz]
Determine el dominio de
\[
f(x)=\sqrt{\frac{x-1}{x+2}}.
\]
\textbf{Solución.} Necesitamos
\[
\frac{x-1}{x+2}\ge0,
\qquad x\neq-2.
\]
Los puntos críticos son $-2$ y $1$. El análisis de signos da
\[
\frac{x-1}{x+2}\ge0
\quad\text{en}\quad
(-\infty,-2)\cup[1,\infty).
\]
Por tanto,
\[
\boxed{\Dom(f)=(-\infty,-2)\cup[1,\infty)}.
\]


\textbf{Gráfica de apoyo.}\medskip
\begin{center}
\begin{tikzpicture}
\begin{axis}[
    width=0.78\textwidth,height=5.4cm,
    axis lines=middle,xmin=-8,xmax=8,ymin=-1,ymax=5,
    xlabel={$x$},ylabel={$y$},grid=both,
    minor tick num=1,samples=180,
    legend style={font=\scriptsize,at={(0.02,0.98)},anchor=north west}
]
\addplot[thick,domain=-8:-2.05] {sqrt((x-1)/(x+2))};
\addplot[thick,domain=1:8] {sqrt((x-1)/(x+2))};
\addplot[dashed] coordinates {(-2,-1) (-2,5)};
\addplot[only marks,mark=*] coordinates {(1,0)};
\end{axis}
\end{tikzpicture}
\end{center}

\end{ejemplo}

\begin{ejemplo}[Reto: raíz en el denominador]
Determine el dominio de
\[
f(x)=\frac{1}{\sqrt{x^2-9}}.
\]
\textbf{Solución.} Aquí no basta pedir que el radicando sea no negativo: la raíz está en el denominador y, por tanto, tampoco puede valer cero. Exigimos
\[
x^2-9>0.
\]
Entonces
\[
|x|>3,
\]
y finalmente
\[
\boxed{\Dom(f)=(-\infty,-3)\cup(3,\infty)}.
\]


\textbf{Gráfica de apoyo.}\medskip
\begin{center}
\begin{tikzpicture}
\begin{axis}[
    width=0.78\textwidth,height=5.4cm,
    axis lines=middle,xmin=-8,xmax=8,ymin=-0.5,ymax=4,
    xlabel={$x$},ylabel={$y$},grid=both,
    minor tick num=1,samples=180,
    legend style={font=\scriptsize,at={(0.02,0.98)},anchor=north west}
]
\addplot[thick,domain=-8:-3.05] {1/sqrt(x^2-9)};
\addplot[thick,domain=3.05:8] {1/sqrt(x^2-9)};
\addplot[dashed] coordinates {(-3,-0.5) (-3,4)};
\addplot[dashed] coordinates {(3,-0.5) (3,4)};
\end{axis}
\end{tikzpicture}
\end{center}

\end{ejemplo}

\begin{ejemplo}[Reto: restricciones simultáneas]
Determine el dominio de
\[
f(x)=\frac{\sqrt{x-1}}{x-4}.
\]
\textbf{Solución.} Deben cumplirse simultáneamente
\[
x-1\ge0
\qquad\text{y}\qquad
x-4\neq0.
\]
Así,
\[
\boxed{\Dom(f)=[1,4)\cup(4,\infty)}.
\]


\textbf{Gráfica de apoyo.}\medskip
\begin{center}
\begin{tikzpicture}
\begin{axis}[
    width=0.78\textwidth,height=5.4cm,
    axis lines=middle,xmin=0,xmax=10,ymin=-6,ymax=6,
    xlabel={$x$},ylabel={$y$},grid=both,
    minor tick num=1,samples=180,
    legend style={font=\scriptsize,at={(0.02,0.98)},anchor=north west}
]
\addplot[thick,domain=1:3.95] {sqrt(x-1)/(x-4)};
\addplot[thick,domain=4.05:10] {sqrt(x-1)/(x-4)};
\addplot[dashed] coordinates {(4,-6) (4,6)};
\addplot[only marks,mark=*] coordinates {(1,0)};
\end{axis}
\end{tikzpicture}
\end{center}

\end{ejemplo}

\subsection{Valor absoluto y funciones definidas por partes}

\begin{ejemplo}[Básico: escribir por partes]
Escriba $f(x)=|x-2|$ como función definida por partes.

\textbf{Solución.} El cambio de signo ocurre cuando $x-2=0$, es decir, en $x=2$. Por definición,
\[
\boxed{
|x-2|=
\begin{cases}
2-x,&x<2,\\
x-2,&x\ge2.
\end{cases}}
\]
El vértice de la gráfica es $(2,0)$.


\textbf{Gráfica de apoyo.}\medskip
\begin{center}
\begin{tikzpicture}
\begin{axis}[
    width=0.78\textwidth,height=5.4cm,
    axis lines=middle,xmin=-3,xmax=7,ymin=-1,ymax=6,
    xlabel={$x$},ylabel={$y$},grid=both,
    minor tick num=1,samples=180,
    legend style={font=\scriptsize,at={(0.02,0.98)},anchor=north west}
]
\addplot[thick,domain=-3:2] {2-x};
\addplot[thick,domain=2:7] {x-2};
\addplot[only marks,mark=*] coordinates {(2,0)};
\end{axis}
\end{tikzpicture}
\end{center}

\end{ejemplo}

\begin{ejemplo}[Intermedio: ecuación con valor absoluto]
Resuelva $|2x-3|=5$.

\textbf{Solución.} La igualdad $|u|=5$ equivale a $u=5$ o $u=-5$:
\[
2x-3=5\Longrightarrow x=4,
\qquad
2x-3=-5\Longrightarrow x=-1.
\]
Por tanto, $\boxed{x\in\{-1,4\}}$.


\textbf{Gráfica de apoyo.}\medskip
\begin{center}
\begin{tikzpicture}
\begin{axis}[
    width=0.78\textwidth,height=5.4cm,
    axis lines=middle,xmin=-3,xmax=6,ymin=-1,ymax=10,
    xlabel={$x$},ylabel={$y$},grid=both,
    minor tick num=1,samples=180,
    legend style={font=\scriptsize,at={(0.02,0.98)},anchor=north west}
]
\addplot[thick,domain=-3:6] {abs(2*x-3)};
\addplot[thick,dashed,domain=-3:6] {5};
\addplot[only marks,mark=*] coordinates {(-1,5) (4,5)};
\end{axis}
\end{tikzpicture}
\end{center}

\end{ejemplo}

\begin{ejemplo}[Intermedio: evaluación por partes]
Sea
\[
f(x)=\begin{cases}
x^2,&x<0,\\
2x+1,&0\le x<3,\\
7,&x\ge3.
\end{cases}
\]
Calcule $f(-2)$, $f(0)$, $f(2)$ y $f(3)$.

\textbf{Solución.} Cada entrada debe evaluarse en la rama que le corresponde:
\[
f(-2)=4,\qquad f(0)=1,\qquad f(2)=5,\qquad f(3)=7.
\]
Los extremos de los intervalos son parte de la definición y no deben ignorarse.


\textbf{Gráfica de apoyo.}\medskip
\begin{center}
\begin{tikzpicture}
\begin{axis}[
    width=0.78\textwidth,height=5.4cm,
    axis lines=middle,xmin=-3,xmax=6,ymin=-1,ymax=10,
    xlabel={$x$},ylabel={$y$},grid=both,
    minor tick num=1,samples=180,
    legend style={font=\scriptsize,at={(0.02,0.98)},anchor=north west}
]
\addplot[thick,domain=-3:-0.02] {x^2};
\addplot[thick,domain=0:2.98] {2*x+1};
\addplot[thick,domain=3:6] {7};
\addplot[only marks,mark=o,mark options={fill=white}] coordinates {(0,0) (3,7)};
\addplot[only marks,mark=*] coordinates {(0,1) (3,7)};
\end{axis}
\end{tikzpicture}
\end{center}

\end{ejemplo}

\subsection{Funciones exponenciales y logarítmicas}

\begin{ejemplo}[Básico: exponencial]
Sea $f(x)=2^x$. Calcule $f(-2)$, $f(0)$ y $f(3)$ e identifique dominio e imagen.

\textbf{Solución.}
\[
f(-2)=\frac14,\qquad f(0)=1,\qquad f(3)=8.
\]
Además,
\[
\boxed{\Dom(f)=\R,\qquad \Ima(f)=(0,\infty)}.
\]


\textbf{Gráfica de apoyo.}\medskip
\begin{center}
\begin{tikzpicture}
\begin{axis}[
    width=0.78\textwidth,height=5.4cm,
    axis lines=middle,xmin=-4,xmax=4,ymin=-1,ymax=10,
    xlabel={$x$},ylabel={$y$},grid=both,
    minor tick num=1,samples=180,
    legend style={font=\scriptsize,at={(0.02,0.98)},anchor=north west}
]
\addplot[thick,domain=-4:4] {2^x};
\addplot[only marks,mark=*] coordinates {(-2,0.25) (0,1) (3,8)};
\end{axis}
\end{tikzpicture}
\end{center}

\end{ejemplo}

\begin{ejemplo}[Intermedio: dominio logarítmico]
Determine el dominio de $f(x)=\ln(5-2x)$.

\textbf{Solución.} El argumento del logaritmo debe ser estrictamente positivo:
\[
5-2x>0\Longrightarrow x<\frac52.
\]
Por tanto,
\[
\boxed{\Dom(f)=(-\infty,5/2)}.
\]


\textbf{Gráfica de apoyo.}\medskip
\begin{center}
\begin{tikzpicture}
\begin{axis}[
    width=0.78\textwidth,height=5.4cm,
    axis lines=middle,xmin=-5,xmax=5,ymin=-6,ymax=5,
    xlabel={$x$},ylabel={$y$},grid=both,
    minor tick num=1,samples=180,
    legend style={font=\scriptsize,at={(0.02,0.98)},anchor=north west}
]
\addplot[thick,domain=-5:2.48] {ln(5-2*x)};
\addplot[dashed] coordinates {(2.5,-6) (2.5,5)};
\end{axis}
\end{tikzpicture}
\end{center}

\end{ejemplo}

\begin{ejemplo}[Reto: logaritmo de una función racional]
Determine el dominio de
\[
f(x)=\ln\left(\frac{x-1}{x+3}\right).
\]
\textbf{Solución.} Exigimos
\[
\frac{x-1}{x+3}>0.
\]
Los puntos críticos son $-3$ y $1$. El cociente es positivo cuando numerador y denominador tienen el mismo signo, por lo que
\[
\boxed{\Dom(f)=(-\infty,-3)\cup(1,\infty)}.
\]


\textbf{Gráfica de apoyo.}\medskip
\begin{center}
\begin{tikzpicture}
\begin{axis}[
    width=0.78\textwidth,height=5.4cm,
    axis lines=middle,xmin=-9,xmax=9,ymin=-5,ymax=5,
    xlabel={$x$},ylabel={$y$},grid=both,
    minor tick num=1,samples=180,
    legend style={font=\scriptsize,at={(0.02,0.98)},anchor=north west}
]
\addplot[thick,domain=-9:-3.05] {ln((x-1)/(x+3))};
\addplot[thick,domain=1.05:9] {ln((x-1)/(x+3))};
\addplot[dashed] coordinates {(-3,-5) (-3,5)};
\addplot[dashed] coordinates {(1,-5) (1,5)};
\end{axis}
\end{tikzpicture}
\end{center}

\end{ejemplo}

\subsection{Transformaciones de gráficas}

\begin{ejemplo}[Intermedio: varias transformaciones]
Describa cómo obtener la gráfica de
\[
g(x)=-2(x-3)^2+1
\]
a partir de $f(x)=x^2$.

\textbf{Solución.} Partiendo de $y=x^2$:
\begin{enumerate}
\item $x\mapsto x-3$: traslación $3$ unidades a la derecha;
\item multiplicar por $2$: estiramiento vertical de factor $2$;
\item signo negativo: reflexión respecto del eje $x$;
\item sumar $1$: traslación una unidad hacia arriba.
\end{enumerate}
El vértice queda en $(3,1)$ y la parábola abre hacia abajo.


\textbf{Gráfica de apoyo.}\medskip
\begin{center}
\begin{tikzpicture}
\begin{axis}[
    width=0.78\textwidth,height=5.4cm,
    axis lines=middle,xmin=-2,xmax=6,ymin=-12,ymax=10,
    xlabel={$x$},ylabel={$y$},grid=both,
    minor tick num=1,samples=180,
    legend style={font=\scriptsize,at={(0.02,0.98)},anchor=north west}
]
\addplot[thick,dashed,domain=-2:6] {x^2};
\addplot[thick,domain=0:6] {-2*(x-3)^2+1};
\addplot[only marks,mark=*] coordinates {(3,1)};
\end{axis}
\end{tikzpicture}
\end{center}

\end{ejemplo}

\begin{ejemplo}[Reto: transformación de una raíz]
Describa la gráfica de $g(x)=\sqrt{4-x}+2$ a partir de $f(x)=\sqrt{x}$ y determine su dominio.

\textbf{Solución.} Como $4-x=-(x-4)$, aparece una reflexión horizontal seguida de una traslación. Para el dominio,
\[
4-x\ge0\Longrightarrow x\le4.
\]
Por tanto,
\[
\boxed{\Dom(g)=(-\infty,4]}.
\]
Además $g(x)\ge2$, de modo que $\Ima(g)=[2,\infty)$.


\textbf{Gráfica de apoyo.}\medskip
\begin{center}
\begin{tikzpicture}
\begin{axis}[
    width=0.78\textwidth,height=5.4cm,
    axis lines=middle,xmin=-5,xmax=8,ymin=-1,ymax=6,
    xlabel={$x$},ylabel={$y$},grid=both,
    minor tick num=1,samples=180,
    legend style={font=\scriptsize,at={(0.02,0.98)},anchor=north west}
]
\addplot[thick,dashed,domain=0:8] {sqrt(x)};
\addplot[thick,domain=-5:4] {sqrt(4-x)+2};
\addplot[only marks,mark=*] coordinates {(4,2)};
\end{axis}
\end{tikzpicture}
\end{center}

\end{ejemplo}

\subsection{Composición e inversas}

\begin{ejemplo}[Intermedio: composición y dominio]
Sean
\[
f(x)=\sqrt{x},\qquad g(x)=x-4.
\]
Determine $f\circ g$ y su dominio.

\textbf{Solución.}
\[
(f\circ g)(x)=f(x-4)=\sqrt{x-4}.
\]
Para que la composición esté definida necesitamos
\[
x-4\ge0,
\]
por lo que
\[
\boxed{\Dom(f\circ g)=[4,\infty)}.
\]
El dominio de una composición no se obtiene simplemente intersectando mecánicamente los dominios: debemos exigir que $g(x)$ pertenezca al dominio de $f$.


\textbf{Gráfica de apoyo.}\medskip
\begin{center}
\begin{tikzpicture}
\begin{axis}[
    width=0.78\textwidth,height=5.4cm,
    axis lines=middle,xmin=0,xmax=12,ymin=-1,ymax=4,
    xlabel={$x$},ylabel={$y$},grid=both,
    minor tick num=1,samples=180,
    legend style={font=\scriptsize,at={(0.02,0.98)},anchor=north west}
]
\addplot[thick,domain=4:12] {sqrt(x-4)};
\addplot[only marks,mark=*] coordinates {(4,0)};
\end{axis}
\end{tikzpicture}
\end{center}

\end{ejemplo}

\begin{ejemplo}[Intermedio: inversa de una función afín]
Determine la inversa de $f(x)=3x-7$.

\textbf{Solución.} Escribimos
\[
y=3x-7
\]
y despejamos $x$:
\[
x=\frac{y+7}{3}.
\]
Intercambiando el nombre de las variables,
\[
\boxed{f^{-1}(x)=\frac{x+7}{3}}.
\]
La comprobación es
\[
f(f^{-1}(x))=3\left(\frac{x+7}{3}\right)-7=x.
\]


\textbf{Gráfica de apoyo.}\medskip
\begin{center}
\begin{tikzpicture}
\begin{axis}[
    width=0.78\textwidth,height=5.4cm,
    axis lines=middle,xmin=-3,xmax=5,ymin=-12,ymax=8,
    xlabel={$x$},ylabel={$y$},grid=both,
    minor tick num=1,samples=180,
    legend style={font=\scriptsize,at={(0.02,0.98)},anchor=north west}
]
\addplot[thick,domain=-3:3] {3*x-7};
\addplot[thick,domain=-3:5] {(x+7)/3};
\addplot[dashed,domain=-3:5] {x};
\end{axis}
\end{tikzpicture}
\end{center}

\end{ejemplo}

\begin{ejemplo}[Reto: restricción necesaria para invertir]
Considere $f(x)=x^2$. Explique por qué no posee inversa en todo $\R$ y encuentre una inversa si restringimos el dominio a $[0,\infty)$.

\textbf{Solución.} En $\R$ la función no es inyectiva, pues $f(a)=f(-a)$. Si restringimos el dominio a $[0,\infty)$, la función sí es inyectiva y su imagen es $[0,\infty)$. De
\[
y=x^2,\qquad x\ge0,
\]
obtenemos
\[
x=\sqrt{y}.
\]
Así,
\[
\boxed{f^{-1}(x)=\sqrt{x},\qquad x\ge0}.
\]


\textbf{Gráfica de apoyo.}\medskip
\begin{center}
\begin{tikzpicture}
\begin{axis}[
    width=0.78\textwidth,height=5.4cm,
    axis lines=middle,xmin=0,xmax=16,ymin=-1,ymax=17,
    xlabel={$x$},ylabel={$y$},grid=both,
    minor tick num=1,samples=180,
    legend style={font=\scriptsize,at={(0.02,0.98)},anchor=north west}
]
\addplot[thick,domain=0:4] {x^2};
\addplot[thick,domain=0:16] {sqrt(x)};
\addplot[dashed,domain=0:8] {x};
\end{axis}
\end{tikzpicture}
\end{center}

\end{ejemplo}

\subsection{Guía rápida de selección para el profesor}

Una posible selección según el desempeño del grupo es:
\begin{center}
\begin{tabular}{p{0.18\textwidth}p{0.23\textwidth}p{0.23\textwidth}p{0.25\textwidth}}
\toprule
\textbf{Familia} & \textbf{Básico} & \textbf{Intermedio} & \textbf{Reto}\\
\midrule
Afines & evaluación & dos puntos & modelación\\
Polinomiales & ceros & vértice/imagen & signo y simetría\\
Racionales & una exclusión & hueco / imagen & varias exclusiones\\
Radicales & desigualdad lineal & cuadrática o cociente & raíz en denominador\\
Valor absoluto & definición por partes & ecuación & análisis gráfico\\
Exp./log. & evaluación & dominio logarítmico & cociente dentro de log\\
Transformaciones & una transformación & varias transformaciones & dominio + transformación\\
Composición/inversa & composición directa & dominio de composición & restricción para invertir\\
\bottomrule
\end{tabular}
\end{center}

La intención de este banco no es aumentar mecánicamente el número de ejercicios, sino disponer de ejemplos con distintos niveles de exigencia para adaptar la explicación al grupo. En cada caso conviene detenerse en tres preguntas: \emph{¿qué información puedo obtener antes de calcular?}, \emph{¿qué restricciones debo respetar?} y \emph{¿qué significa geométricamente el resultado?}


\section{Ejemplos integradores}

\begin{ejemplo}[Análisis de una función racional]
Considere
\[
f(x)=\frac{x+1}{x-2}.
\]

\textbf{Dominio.}
\[
\Dom(f)=\R\setminus\{2\}.
\]

\textbf{Intersección con el eje $x$.}
\[
f(x)=0
\Longleftrightarrow x+1=0
\Longleftrightarrow x=-1.
\]

\textbf{Intersección con el eje $y$.}
\[
f(0)=\frac{1}{-2}=-\frac12.
\]

\textbf{Imagen.}
Sea
\[
y=\frac{x+1}{x-2}.
\]
Despejando $x$:
\[
yx-2y=x+1,
\]
\[
x(y-1)=2y+1.
\]
Si $y\neq1$,
\[
x=\frac{2y+1}{y-1}.
\]
Por tanto, todos los reales excepto $y=1$ son alcanzados:
\[
\Ima(f)=\R\setminus\{1\}.
\]
\end{ejemplo}

\begin{ejemplo}[Composición y restricciones]
Sean
\[
f(x)=\frac{1}{x-1},
\qquad
g(x)=\sqrt{x+2}.
\]
Entonces
\[
(f\circ g)(x)
=\frac{1}{\sqrt{x+2}-1}.
\]
Se requiere
\[
x+2\geq0
\]
y, además,
\[
\sqrt{x+2}\neq1.
\]
La primera condición da $x\geq-2$ y la segunda
\[
x+2\neq1\Longleftrightarrow x\neq-1.
\]
Así,
\[
\Dom(f\circ g)=[-2,-1)\cup(-1,\infty).
\]
\end{ejemplo}

% =========================================================
\section{Errores frecuentes}

\begin{enumerate}
    \item Confundir codominio con imagen.
    \item Pensar que una función no puede asignar el mismo valor a dos entradas.
    \item Confundir $f(x+h)$ con $f(x)+h$.
    \item Cancelar factores sin conservar las restricciones del dominio original.
    \item Usar $\sqrt{g(x)}$ sin imponer $g(x)\geq0$.
    \item Usar $\ln(g(x))$ imponiendo $g(x)\geq0$ en vez de $g(x)>0$.
    \item Suponer que $\Dom(f\circ g)$ es simplemente
    $\Dom(f)\cap\Dom(g)$.
    \item Confundir $f^{-1}(x)$ con $1/f(x)$.
    \item Despejar una ``inversa'' sin verificar que la función sea uno a uno en
    el dominio considerado.
    \item Ignorar las restricciones físicas al construir un modelo.
\end{enumerate}
\newpage
% =========================================================
\section{Ejercicios}

\subsection{Concepto de función}

Determine cuáles de las siguientes relaciones definen una función de $A$ en
$B$. Justifique formalmente su respuesta.

\begin{enumerate}
    \item $A=\{1,2,3\}$, $B=\{a,b\}$,
    \[
    R=\{(1,a),(2,a),(3,b)\}.
    \]
    \item
    \[
    R=\{(1,a),(1,b),(2,a),(3,b)\}.
    \]
    \item La relación en $\R$ dada por
    \[
    y^2=x.
    \]
    \item La relación
    \[
    y=x^3-2x+1.
    \]
\end{enumerate}

\subsection{Evaluación}

Sea
\[
f(x)=2x^2-5x+3.
\]
Calcule:
\begin{multicols}{2}
\begin{enumerate}
    \item $f(0)$.
    \item $f(2)$.
    \item $f(-3)$.
    \item $f(a)$.
    \item $f(a+1)$.
    \item $f(x+h)$.
    \item $f(x+h)-f(x)$.
    \item $\dfrac{f(x+h)-f(x)}{h}$, $h\neq0$.
\end{enumerate}
\end{multicols}

\subsection{Dominio}

Determine el dominio de cada función y expréselo en notación de intervalos.

\begin{enumerate}
    \item $f(x)=3x^4-2x+7$.
    \item $f(x)=\dfrac{x+3}{x-5}$.
    \item $f(x)=\dfrac{x^2-1}{x^2-9}$.
    \item $f(x)=\sqrt{7-2x}$.
    \item $f(x)=\sqrt{x^2-9}$.
    \item $f(x)=\dfrac{1}{\sqrt{x-2}}$.
    \item $f(x)=\ln(x+4)$.
    \item $f(x)=\ln(9-x^2)$.
    \item $f(x)=\sqrt{\dfrac{x-2}{x+1}}$.
    \item $f(x)=\dfrac{\sqrt{2x+1}}{x^2-4}$.
    \item $f(x)=\ln\left(\dfrac{x-1}{x+3}\right)$.
    \item $f(x)=\sqrt{\ln x}$.
\end{enumerate}

\subsection{Simetría}

Determine algebraicamente si cada función es par, impar o ninguna.

\begin{enumerate}
    \item $f(x)=x^6+3x^2-1$.
    \item $f(x)=x^5-4x^3+x$.
    \item $f(x)=x^3+x^2$.
    \item $f(x)=\dfrac{x}{x^2+1}$.
    \item $f(x)=\dfrac{1}{x^2+4}$.
\end{enumerate}

\subsection{Operaciones}

Sean
\[
f(x)=\sqrt{x+1},
\qquad
g(x)=\frac{1}{x-2}.
\]
Obtenga las siguientes funciones y determine cuidadosamente sus dominios:
\[
f+g,\qquad f-g,\qquad fg,\qquad\frac fg,\qquad\frac gf.
\]

\subsection{Composición}

Para cada par, calcule $f\circ g$ y $g\circ f$ y determine sus dominios.

\begin{enumerate}
    \item
    \[
    f(x)=x^2+1,\qquad g(x)=2x-3.
    \]
    \item
    \[
    f(x)=\sqrt{x},\qquad g(x)=x-4.
    \]
    \item
    \[
    f(x)=\frac1x,\qquad g(x)=x+2.
    \]
    \item
    \[
    f(x)=\ln x,\qquad g(x)=x^2-1.
    \]
\end{enumerate}

\subsection{Inyectividad e inversas}

\begin{enumerate}
    \item Demuestre que
    \[
    f:\R\to\R,\qquad f(x)=5x+2
    \]
    es biyectiva y determine $f^{-1}$.

    \item Determine por qué
    \[
    f:\R\to\R,\qquad f(x)=x^2
    \]
    no tiene inversa.

    \item Considere ahora
    \[
    f:[0,\infty)\to[0,\infty),\qquad f(x)=x^2.
    \]
    Demuestre que es biyectiva y obtenga su inversa.

    \item Determine la inversa de
    \[
    f(x)=\frac{2x-1}{x+3},
    \]
    indicando apropiadamente dominio y codominio.
\end{enumerate}

\subsection{Modelación}

En cada problema:
\begin{enumerate}[label=\alph*)]
    \item identifique las variables;
    \item elija una variable independiente;
    \item construya la función solicitada;
    \item determine el dominio físicamente admisible.
\end{enumerate}

\begin{enumerate}
    \item Un rectángulo tiene perímetro $60$ cm. Exprese su área como función
    de uno de sus lados.

    \item Se dispone de $100$ m de cerca para construir un corral rectangular
    junto a un muro, por lo que sólo deben cercarse tres lados. Exprese el área
    como función del ancho perpendicular al muro.

    \item Una caja sin tapa se construye cortando cuadrados de lado $x$ en las
    esquinas de una lámina rectangular de $30$ cm por $20$ cm y doblando los
    bordes. Exprese el volumen como función de $x$.

    \item Una compañía cobra \$80 de tarifa fija y \$12 por kilómetro recorrido.
    Construya la función de costo.

    \item El precio de un artículo se modela por
    \[
    p(x)=600-3x,
    \]
    donde $x$ es la cantidad vendida. Determine la función de ingreso.

    \item Si en el problema anterior el costo es
    \[
    C(x)=5000+120x,
    \]
    determine la función de utilidad.

    \item Una población inicial de $5000$ individuos crece $3\%$ por año.
    Construya un modelo exponencial discreto.

    \item Un automóvil debe recorrer $300$ km. Exprese el tiempo de viaje como
    función de la velocidad promedio.
\end{enumerate}



\section*{Ejercicios 1.1}

\textit{En los ejercicios 1 a 4, determine si el conjunto es una función. Si es una función determine su dominio.}

\begin{enumerate}
\item
\begin{enumerate}[label=(\alph*)]
\item $\{(x,y)\mid y=\sqrt{x-4}\}$
\item $\{(x,y)\mid y=\sqrt{x^2-4}\}$
\item $\{(x,y)\mid y=\sqrt{4-x^2}\}$
\item $\{(x,y)\mid x^2+y^2=4\}$
\end{enumerate}

\item
\begin{enumerate}[label=(\alph*)]
\item $\{(x,y)\mid y=\sqrt{x+1}\}$
\item $\{(x,y)\mid y=\sqrt{x^2-1}\}$
\item $\{(x,y)\mid y=\sqrt{1-x^2}\}$
\item $\{(x,y)\mid x^2+y^2=1\}$
\end{enumerate}

\item
\begin{enumerate}[label=(\alph*)]
\item $\{(x,y)\mid y=x^2\}$
\item $\{(x,y)\mid x=y^2\}$
\item $\{(x,y)\mid y=x^3\}$
\item $\{(x,y)\mid x=y^3\}$
\end{enumerate}

\item
\begin{enumerate}[label=(\alph*)]
\item $\{(x,y)\mid y=(x-1)^2+2\}$
\item $\{(x,y)\mid x=(y-2)^2+1\}$
\item $\{(x,y)\mid y=(x+2)^3-1\}$
\item $\{(x,y)\mid x=(y+1)^3-2\}$
\end{enumerate}

\item Dada $f(x)=2x-1$, determine
\begin{enumerate}[label=(\alph*)]
\item $f(3)$; \item $f(-2)$; \item $f(0)$; \item $f(a+1)$; \item $f(x+1)$;
\item $f(2x)$; \item $2f(x)$; \item $f(x+h)$; \item $f(x)+f(h)$;
\item $\dfrac{f(x+h)-f(x)}{h}$, $h\ne0$.
\end{enumerate}

\item Dada $f(x)=\dfrac{3}{x}$, calcule
\begin{enumerate}[label=(\alph*)]
\item $f(1)$; \item $f(-3)$; \item $f(6)$; \item $f\!\left(\dfrac13\right)$;
\item $f\!\left(\dfrac3a\right)$; \item $f\!\left(\dfrac3x\right)$;
\item $\dfrac{f(3)}{f(x)}$; \item $f(x-3)$; \item $f(x)-f(3)$;
\item $\dfrac{f(x+h)-f(x)}{h}$, $h\ne0$.
\end{enumerate}

\item Dada $f(x)=2x^2+5x-3$, determine
\begin{enumerate}[label=(\alph*)]
\item $f(-2)$; \item $f(-1)$; \item $f(0)$; \item $f(3)$; \item $f(h+1)$;
\item $f(2x^2)$; \item $f(x^2-3)$; \item $f(x+h)$; \item $f(x)+f(h)$;
\item $\dfrac{f(x+h)-f(x)}{h}$, $h\ne0$.
\end{enumerate}

\item Dada $g(x)=3x^2-4$, calcule
\begin{enumerate}[label=(\alph*)]
\item $g(-4)$; \item $g\!\left(\dfrac12\right)$; \item $g(x^2)$; \item $g(3x^2-4)$;
\item $g(x-h)$; \item $g(x)-g(h)$;
\item $\dfrac{g(x+h)-g(x)}{h}$, $h\ne0$.
\end{enumerate}

\item Dada $F(x)=\sqrt{x+9}$, encuentre
\begin{enumerate}[label=(\alph*)]
\item $F(x+9)$; \item $F(x^2-9)$; \item $F(x^4-9)$; \item $F(x^2+6x)$;
\item $F(x^4-6x^2)$; \item $\dfrac{F(x+h)-F(x)}{h}$, $h\ne0$.
\end{enumerate}

\item Dada $G(x)=\sqrt{4-x}$, determine
\begin{enumerate}[label=(\alph*)]
\item $G(4-x)$; \item $G(4-x^2)$; \item $G(4-x^4)$; \item $G(4x-x^2)$;
\item $G(-x^4-4x^2)$; \item $\dfrac{G(x+h)-G(x)}{h}$, $h\ne0$.
\end{enumerate}
\end{enumerate}

\textit{En los ejercicios 11 a 46, dibuje a mano la gráfica de la función y determine su dominio y su contradominio.}

\begin{enumerate}[resume]
\item $f(x)=3x-1$
\item $g(x)=4-x$
\item $F(x)=2x^2$
\item $G(x)=x^2+2$
\item $g(x)=5-x^2$
\item $f(x)=(x-1)^2$
\item $G(x)=\sqrt{x-1}$
\item $F(x)=\sqrt{9-x}$
\item $f(x)=\sqrt{x^2-4}$
\item $g(x)=\sqrt{4-x^2}$
\item $g(x)=\sqrt{9-x^2}$
\item $f(x)=\sqrt{x^2-1}$
\item $h(x)=|x-3|$
\item $H(x)=|5-x|$
\item $F(x)=|3x+2|$
\item $G(x)=\dfrac{x^2-4}{x-2}$
\item $H(x)=\dfrac{x^2-25}{x+5}$
\item $f(x)=\dfrac{2x^2+7x+3}{x+3}$
\item $f(x)=\dfrac{x^2-4x+3}{x-1}$
\item $g(x)=\dfrac{(x^2-4)(x-3)}{x^2-x-6}$
\item $f(x)=\begin{cases}-2,&x\le3,\\2,&3<x.\end{cases}$
\item $g(x)=\begin{cases}-4,&x<-2,\\-1,&-2\le x\le2,\\3,&2<x.\end{cases}$
\item $g(x)=\begin{cases}2x-1,&x\ne2,\\0,&x=2.\end{cases}$
\item $f(x)=\begin{cases}3x+2,&x\ne1,\\8,&x=1.\end{cases}$
\item $F(x)=\begin{cases}x^2-4,&x\ne3,\\-2,&x=3.\end{cases}$
\item $G(x)=\begin{cases}9-x^2,&x\ne-3,\\4,&x=-3.\end{cases}$
\item $G(x)=\begin{cases}1-x^2,&x<0,\\3x+1,&0\le x.\end{cases}$
\item $F(x)=\begin{cases}x^2-4,&x<3,\\2x-1,&3\le x.\end{cases}$
\item $g(x)=\begin{cases}6x+7,&x\le-2,\\4-x,&-2<x.\end{cases}$
\item $f(x)=\begin{cases}x-2,&x\le0,\\x^2+1,&0<x.\end{cases}$
\item $h(x)=\begin{cases}x+3,&x<-5,\\\sqrt{25-x^2},&-5\le x\le5,\\3-x,&5<x.\end{cases}$
\item $H(x)=\begin{cases}\dfrac{x+2}{\sqrt{16-x^2}},&x\le-4,\\2-x,&-4<x<4,\\2-x,&4\le x.\end{cases}$
\item $F(x)=\dfrac{x^3-2x^2}{x-2}$
\item $G(x)=\dfrac{x^3+3x^2}{x+3}$
\item $f(x)=\lfloor x-4\rfloor$
\item $g(x)=\lfloor x+2\rfloor$
\end{enumerate}

\begin{enumerate}[start=47]
\item
\begin{enumerate}[label=(\alph*)]
\item Dibuje la gráfica de la función escalón (o salto) unitario denotada por $U$ y definida por
\[
U(x)=\begin{cases}0,&x<0,\\1,&0\le x.\end{cases}
\]
Defina cada una de las siguientes funciones a trozos y dibuje sus gráficas:
\item $U(x-1)$; \item $U(x)-1$; \item $U(x)-U(x-1)$.
\end{enumerate}

\item Defina cada una de las siguientes funciones a trozos y dibuje sus gráficas, donde $U$ es la función escalón unitario definida en el ejercicio 47:
\begin{enumerate}[label=(\alph*)]
\item $x\,U(x)$; \item $(x+1)U(x+1)$; \item $(x+1)U(x+1)-xU(x)$.
\end{enumerate}

\item
\begin{enumerate}[label=(\alph*)]
\item Dibuje la gráfica de la función signo denotada por $\operatorname{sgn}$ y definida por
\[
\operatorname{sgn}x=\begin{cases}-1,&x<0,\\0,&x=0,\\1,&0<x.\end{cases}
\]
$\operatorname{sgn}(x)$ se lee ``signo de $x$''. Defina cada una de las siguientes funciones a trozos y dibuje sus gráficas:
\item $x\,\operatorname{sgn}(x)$; \item $2-x\,\operatorname{sgn}(x)$; \item $x-2\,\operatorname{sgn}(x)$.
\end{enumerate}

\item Defina cada una de las siguientes funciones a trozos, donde $\operatorname{sgn}$ es la función signo definida en el ejercicio 49:
\begin{enumerate}[label=(\alph*)]
\item $\operatorname{sgn}(x+1)$; \item $\operatorname{sgn}(x-1)$; \item $\operatorname{sgn}(x+1)-\operatorname{sgn}(x-1)$.
\end{enumerate}

\item La gráfica de la función $f$ de la figura se parece a la letra W. Defina $f(x)$ a trozos.
\item La gráfica de la función $f$ de la figura se parece a la letra M. Defina $f(x)$ a trozos.
\item En la figura, la gráfica se parece a la letra X y es la gráfica de dos funciones $f_1$ y $f_2$ trazadas en el rectángulo de inspección de $[-1,1]$ por $[-1,1]$. Defina $f_1(x)$ y $f_2(x)$.
\item Existen tres funciones $f_1$, $f_2$ y $f_3$ cuyas gráficas trazadas simultáneamente en el rectángulo de inspección de $[-1,1]$ por $[-1,1]$ se parecen a la letra Z. Defina $f_1(x)$, $f_2(x)$ y $f_3(x)$.
\end{enumerate}

\textit{En los ejercicios 55 a 58, haga lo siguiente: (a) defina la función a trozos sin emplear las barras de valor absoluto; (b) dibuje la gráfica de la función definida en el inciso (a); (c) apoye sus respuestas a los incisos (a) y (b) trazando la gráfica de la función.}

\begin{enumerate}[start=55]
\item $f(x)=|x^2-1|$
\item $g(x)=|4-x^2|$
\item $g(x)=|x|\,|5-x|$
\item $f(x)=|x|\,|x-3|$
\end{enumerate}

\newpage
\section*{Ejercicios 1.2}

\textit{En los ejercicios 1 a 10, defina las siguientes funciones y determine el dominio de la función resultante: (a) $f+g$; (b) $f-g$; (c) $f\cdot g$; (d) $f/g$; (e) $g/f$.}

\begin{enumerate}
\item $f(x)=x-5$; $g(x)=x^2-1$
\item $f(x)=\sqrt{x}$; $g(x)=x^2+1$
\item $f(x)=\dfrac{x+1}{x-1}$; $g(x)=\dfrac1x$
\item $f(x)=\sqrt{x}$; $g(x)=4-x^2$
\item $f(x)=\sqrt{x}$; $g(x)=x^2-1$
\item $f(x)=|x|$; $g(x)=|x-3|$
\item $f(x)=x^2+1$; $g(x)=3x-2$
\item $f(x)=\sqrt{x-4}$; $g(x)=x^2-4$
\item $f(x)=\dfrac1{x+1}$; $g(x)=\dfrac{x}{x-2}$
\item $f(x)=x^2$; $g(x)=\dfrac1{\sqrt{x}}$
\end{enumerate}

\textit{En los ejercicios 11 a 14, para las funciones $f$ y $g$ y el número $c$, obtenga $(f\circ g)(c)$ mediante dos métodos: (a) calcule $g(c)$ y utilice este número para determinar $f(g(c))$; (b) determine $(f\circ g)(x)$ y emplee este valor para calcular $(f\circ g)(c)$.}

\begin{enumerate}[start=11]
\item $f(x)=3x^2-4x$; $g(x)=2x-5$; $c=4$
\item $f(x)=\sqrt{x^2-36}$; $g(x)=x^2-3x$; $c=5$
\item $f(x)=\dfrac1{x-1}$; $g(x)=\dfrac2{x^2+1}$; $c=\dfrac12$
\item $f(x)=\dfrac{2\sqrt{x}+3}{x}$; $g(x)=\dfrac{2x+5}{x^4}$; $c=-2$
\end{enumerate}

\textit{En los ejercicios 15 a 24, defina las siguientes funciones y determine el dominio de la función compuesta: (a) $f\circ g$; (b) $g\circ f$; (c) $f\circ f$; (d) $g\circ g$.}

\begin{enumerate}[start=15]
\item $f(x)=x-2$; $g(x)=x+7$
\item $f(x)=3-2x$; $g(x)=6-3x$
\item Las funciones del ejercicio 1.
\item Las funciones del ejercicio 2.
\item $f(x)=\sqrt{x+2}$; $g(x)=x^2-2$
\item $f(x)=x^2-1$; $g(x)=\dfrac1x$
\item $f(x)=\dfrac1x$; $g(x)=\sqrt{x}$
\item $f(x)=\sqrt{x}$; $g(x)=-\dfrac1x$
\item $f(x)=|x|$; $g(x)=|x+2|$
\item $f(x)=\sqrt{x^2-1}$; $g(x)=\sqrt{x-1}$
\end{enumerate}

\textit{En los ejercicios 25 y 26 defina las siguientes funciones y determine el dominio de las funciones resultantes: (a) $f(x^2)$; (b) $[f(x)]^2$; (c) $f\circ f(x)$; (d) $f\circ f(-x)$.}

\begin{enumerate}[start=25]
\item $f(x)=\sqrt{x}$
\item $f(x)=\dfrac1{x-1}$
\end{enumerate}

\textit{En los ejercicios 27 a 32, exprese $h$ como composición de las dos funciones $f$ y $g$ en dos formas.}

\begin{enumerate}[start=27]
\item $h(x)=\sqrt{x^2-4}$
\item $h(x)=(9+x^2)^{-2}$
\item $h(x)=\left(\dfrac1{x-2}\right)^3$
\item $h(x)=\dfrac4{\sqrt[3]{x^3+3}}$
\item $h(x)=(x^2+4x-5)^4$
\item $h(x)=\sqrt{|x|+4}$
\end{enumerate}

\textit{En los ejercicios 33 a 38, trace en la graficadora la gráfica de la función y a partir de la gráfica conjeture si la función es par, impar o de ninguno de estos dos tipos. Después confirme su conjetura analíticamente.}

\begin{enumerate}[start=33]
\item \begin{enumerate}[label=(\alph*)]\item $f(x)=2x^4-3x^2+1$ \item $g(x)=5x^5+1$\end{enumerate}
\item \begin{enumerate}[label=(\alph*)]\item $f(x)=x^2+2x+2$ \item $g(x)=x^6-1$\end{enumerate}
\item \begin{enumerate}[label=(\alph*)]\item $f(x)=5x^3-7x$ \item $g(x)=|x|$\end{enumerate}
\item \begin{enumerate}[label=(\alph*)]\item $f(x)=4x^5+3x^3$ \item $g(x)=x^3+1$\end{enumerate}
\item \begin{enumerate}[label=(\alph*)]\item $f(x)=\sqrt[3]{x}$ \item $g(x)=5x^4-4$\end{enumerate}
\item \begin{enumerate}[label=(\alph*)]\item $f(x)=\dfrac{|x|}{x}$ \item $g(x)=2|x|+3$\end{enumerate}
\end{enumerate}

\textit{En los ejercicios 39 y 40, determine analíticamente si la función es par, impar o de ninguno de estos dos tipos.}

\begin{enumerate}[start=39]
\item \begin{enumerate}[label=(\alph*)]\item $f(y)=\dfrac{y^3-y}{y^2+1}$ \item $g(r)=\dfrac{r^2-1}{r^2+1}$ \item $f(x)=\dfrac{|x|}{x^2+1}$\end{enumerate}
\item \begin{enumerate}[label=(\alph*)]\item $h(x)=\dfrac{x^2-5}{2x^3+x}$ \item $g(z)=\dfrac{z-1}{z+1}$ \item $f(x)=\begin{cases}-1,&x<0,\\1,&0\le x.\end{cases}$\end{enumerate}
\end{enumerate}

\textit{En los ejercicios 41 a 44, haga lo siguiente: (a) defina $f(x)$, sin las barras de valor absoluto, en los intervalos indicados. (b) Apoye la respuesta del inciso (a), gráficamente trazando la gráfica de $f$ en la graficadora a partir de la ecuación dada. (c) A partir de la gráfica del inciso (b), establezca si la función es par, impar o de ninguno de estos dos tipos. (d) Confirme la respuesta del inciso (c) analíticamente a partir de la ecuación dada.}

\begin{enumerate}[start=41]
\item $f(x)=\dfrac{|x|}{x}$; $(-\infty,0)$, $(0,+\infty)$
\item $f(x)=x|x|$; $(-\infty,0)$, $[0,+\infty)$
\item $f(x)=|x-2|-|x+2|$; $(-\infty,-2)$, $[-2,2)$, $[2,+\infty)$
\item $f(x)=\dfrac{|x+1|-|x-1|}{x}$; $(-\infty,-1)$, $[-1,0)$, $(0,1]$, $(1,+\infty)$
\end{enumerate}

\begin{enumerate}[start=45]
\item ¿Es conmutativa la composición de dos funciones? Es decir, si $f$ y $g$ son dos funciones cualesquiera, ¿son iguales $(f\circ g)(x)$ y $(g\circ f)(x)$? Justifique su respuesta proporcionando un ejemplo.
\end{enumerate}

Si $f$ y $g$ son dos funciones tales que $(f\circ g)(x)=x$ y $(g\circ f)(x)=x$, entonces se dice que $f$ y $g$ son inversas una de la otra. En los ejercicios 46 a 50, demuestre que $f$ y $g$ son inversas una de la otra.

\begin{enumerate}[start=46]
\item $f(x)=2x-3$ y $g(x)=\dfrac{x+3}{2}$
\item $f(x)=\dfrac1{x+1}$ y $g(x)=\dfrac{1-x}{x}$
\item $f(x)=x^2$, $x\ge0$, y $g(x)=\sqrt{x}$
\item $f(x)=x^2$, $x\le0$, y $g(x)=-\sqrt{x}$
\item $f(x)=(x-1)^3$ y $g(x)=1+\sqrt[3]{x}$
\item La función escalón unitario $U$ y la función signo $\operatorname{sgn}$ se definieron en los ejercicios 47 y 49, respectivamente, de la sección 1.1. (a) Defina $\operatorname{sgn}(U(x))$ y dibuje la gráfica. (b) Defina $U(\operatorname{sgn}(x))$ y dibuje la gráfica.
\item Demuestre que si $f$ y $g$ son funciones impares, entonces $(f+g)$ y $(f-g)$ también son funciones impares, mientras que $fg$ y $f/g$ son funciones pares.
\item Determine si la función compuesta $f\circ g$ es par o impar en cada uno de los casos siguientes: (a) $f$ y $g$ son impares; (b) $f$ es par y $g$ es impar; (c) $g$ es par.
\item Encuentre fórmulas para $(f\circ g)(x)$ si
\[
f(x)=\begin{cases}0,&x<0,\\2x,&0\le x\le1,\\0,&1<x,\end{cases}
\qquad
y\qquad
g(x)=\begin{cases}1,&x<0,\\\frac12x,&0\le x\le1,\\1,&1<x.\end{cases}
\]
Dibuje las gráficas de $f$, $g$ y $f\circ g$.
\item Encuentre fórmulas para $(g\circ f)(x)$ a partir de las funciones del ejercicio 54. Dibuje la gráfica de $g\circ f$.
\item Si $f(x)=x^2+2x+2$, encuentre dos funciones $g$ para las cuales $(f\circ g)(x)=x^2-4x+5$.
\item Si $f(x)=x^2$, encuentre dos funciones $g$ para las cuales $(f\circ g)(x)=4x^2-12x+9$.
\item Demuestre que si $f$ y $g$ son dos funciones lineales, entonces $f\circ g$ es una función lineal.
\item Existe una función cuyo dominio es el conjunto de todos los números reales que es a la vez par e impar. ¿Cuál es esa función? Demuestre que es única esta función.
\item Suponga que $f(x)=\dfrac1x$, $g(x)=-\dfrac1x$ y $h(x)=-x$. Demuestre que $(f\circ g)(x)=(g\circ f)(x)$ y explique por qué $f\circ g$ ni $g\circ f$ son la misma que $h$.
\item Trace en la graficadora las gráficas de las dos funciones $F$ y $G$ definidas por
\[
F(x)=\frac{\sqrt{x+1}}{\sqrt{x-4}}
\qquad\text{y}\qquad
G(x)=\sqrt{\frac{x+1}{x-4}}.
\]
Observe que $F$ es la misma función que $f/g$ del ejemplo 1(d). Explique por qué las gráficas de $F$ y $G$ no son las mismas y, consecuentemente, las funciones no son iguales.
\end{enumerate}


\newpage
% =========================================================
\section{Problemas de profundización}

\begin{enumerate}
    \item Sean $f,g:\R\to\R$. Demuestre que si $f$ y $g$ son pares, entonces
    $f+g$ y $fg$ son pares.

    \item Demuestre que si $f$ y $g$ son impares, entonces $f+g$ es impar,
    mientras que $fg$ es par.

    \item Demuestre que la composición de dos funciones inyectivas es
    inyectiva.

    \item Demuestre que la composición de dos funciones suprayectivas es
    suprayectiva, suponiendo compatibles los dominios y codominios.

    \item Si $f$ y $g$ son biyectivas, demuestre que
    \[
    (f\circ g)^{-1}=g^{-1}\circ f^{-1}.
    \]

    \item Sea
    \[
    f(x)=\frac{ax+b}{cx+d},
    \qquad ad-bc\neq0.
    \]
    Determine su dominio y despeje formalmente $x$ en términos de $y$ para
    obtener la función inversa, especificando las restricciones necesarias.

    \item Determine la imagen de
    \[
    f(x)=x^2-4x+7
    \]
    completando el cuadrado.

    \item Determine la imagen de
    \[
    f(x)=\frac{2x+1}{x-3}
    \]
    despejando $x$ en términos de $y$.
\end{enumerate}

% =========================================================
\section{Resumen conceptual}

Al finalizar esta unidad el estudiante debe distinguir con claridad los
siguientes objetos:

\[
\boxed{
f:A\to B,\qquad x\mapsto f(x)
}
\]

\begin{itemize}
    \item $A$: dominio;
    \item $B$: codominio;
    \item $f(A)$: imagen;
    \item $G_f=\{(x,f(x)):x\in A\}$: gráfica;
    \item $(f\circ g)(x)=f(g(x))$: composición;
    \item $f^{-1}$: inversa, cuando $f$ es biyectiva.
\end{itemize}

La función debe entenderse simultáneamente desde cuatro perspectivas:
\begin{enumerate}
    \item \textbf{conjuntista}: como una correspondencia con existencia y
    unicidad;
    \item \textbf{algebraica}: mediante una regla de correspondencia;
    \item \textbf{geométrica}: mediante su gráfica;
    \item \textbf{aplicada}: como un modelo de dependencia entre cantidades.
\end{enumerate}

Esta visión prepara el terreno para la siguiente unidad. En efecto, el concepto
de límite estudiará el comportamiento de $f(x)$ cuando $x$ se aproxima a un
valor determinado, y la continuidad relacionará dicho comportamiento con el
valor de la propia función.

% =========================================================
\appendix
\section{Apéndice: inyectividad, biyectividad y función inversa}
\label{ap:inyectividad-inversa}
% =========================================================

\subsection{Dominio, codominio e imagen}
Sea $f:A\to B$. El conjunto $A$ es el dominio, $B$ es el codominio y
\[
\operatorname{Im}(f)=\{f(x):x\in A\}\subseteq B
\]
es la imagen o rango. La distinción entre imagen y codominio es esencial para estudiar suprayectividad y existencia de inversa.

\begin{ejemplo}
Para $f:\mathbb R\to\mathbb R$, $f(x)=x^2$, se tiene
\[
\operatorname{Im}(f)=[0,\infty),
\]
de modo que la imagen no coincide con el codominio declarado.
\end{ejemplo}

\subsection{Funciones inyectivas}
\begin{definicion}
Una función $f:A\to B$ es \textbf{inyectiva} si elementos distintos del dominio tienen imágenes distintas. Equivalentemente,
\[
\boxed{f(x_1)=f(x_2)\Longrightarrow x_1=x_2.}
\]
\end{definicion}

\begin{ejemplo}
Sea $f(x)=3x-2$. Si $f(x_1)=f(x_2)$, entonces
\[
3x_1-2=3x_2-2\Longrightarrow 3x_1=3x_2\Longrightarrow x_1=x_2.
\]
Por tanto, $f$ es inyectiva.
\end{ejemplo}

\begin{ejemplo}
La función $f(x)=x^2$ no es inyectiva en $\mathbb R$, pues
\[
f(2)=f(-2)=4,\qquad 2\ne-2.
\]
\end{ejemplo}

\subsubsection{Prueba de la recta horizontal}
Una función es inyectiva si toda recta horizontal $y=c$ corta su gráfica a lo más una vez. No debe confundirse con la prueba de la recta vertical: ésta determina si una gráfica representa una función; la horizontal estudia inyectividad.

\subsection{Funciones suprayectivas}
\begin{definicion}
Una función $f:A\to B$ es \textbf{suprayectiva} si todo elemento del codominio es imagen de al menos un elemento del dominio:
\[
\boxed{\forall y\in B\;\exists x\in A:\ f(x)=y.}
\]
Equivalentemente, $\operatorname{Im}(f)=B$.
\end{definicion}

\begin{ejemplo}
Sea $f:\mathbb R\to\mathbb R$, $f(x)=2x+3$. Dado $y\in\mathbb R$, resolver
\[
y=2x+3
\]
produce
\[
x=\frac{y-3}{2}\in\mathbb R.
\]
Por tanto, todo $y$ es alcanzado y $f$ es suprayectiva.
\end{ejemplo}

\begin{observacion}
La suprayectividad depende del codominio. La regla $f(x)=x^2$ no es suprayectiva como función $\mathbb R\to\mathbb R$, pero sí lo es como función $\mathbb R\to[0,\infty)$.
\end{observacion}

\subsection{Funciones biyectivas}
\begin{definicion}
Una función es \textbf{biyectiva} si es simultáneamente inyectiva y suprayectiva. Así, cada elemento del codominio corresponde a exactamente un elemento del dominio.
\end{definicion}

\begin{ejemplo}
Sea $f:\mathbb R\to\mathbb R$, $f(x)=5x-7$. La igualdad
\[
5x_1-7=5x_2-7
\]
implica $x_1=x_2$, por lo que es inyectiva. Además, dado $y\in\mathbb R$,
\[
y=5x-7\Longrightarrow x=\frac{y+7}{5}\in\mathbb R,
\]
por lo que es suprayectiva. En consecuencia, $f$ es biyectiva.
\end{ejemplo}

\subsection{Función inversa}
\begin{definicion}
Si $f:A\to B$ es biyectiva, su \textbf{función inversa} es la función
\[
f^{-1}:B\to A
\]
que satisface
\[
f^{-1}(y)=x\Longleftrightarrow f(x)=y.
\]
\end{definicion}

Las identidades fundamentales son
\[
\boxed{f^{-1}(f(x))=x}\quad(x\in A),
\qquad
\boxed{f(f^{-1}(y))=y}\quad(y\in B).
\]

\begin{observacion}
No debe confundirse la inversa con el recíproco:
\[
\boxed{f^{-1}(x)\neq \frac{1}{f(x)}}
\]
en general.
\end{observacion}

\subsubsection{Procedimiento algebraico}
Para hallar $f^{-1}$:
\begin{enumerate}
\item escribir $y=f(x)$;
\item intercambiar $x$ y $y$;
\item despejar $y$;
\item escribir $y=f^{-1}(x)$;
\item verificar, cuando convenga, mediante composición.
\end{enumerate}

\begin{ejemplo}
Sea $f(x)=3x-5$. Partimos de
\[
y=3x-5.
\]
Intercambiando variables,
\[
x=3y-5.
\]
Entonces
\[
x+5=3y,\qquad y=\frac{x+5}{3},
\]
y por tanto
\[
\boxed{f^{-1}(x)=\frac{x+5}{3}.}
\]
La verificación es
\[
f\!\left(f^{-1}(x)\right)=3\left(\frac{x+5}{3}\right)-5=x
\]
y
\[
f^{-1}(f(x))=\frac{(3x-5)+5}{3}=x.
\]
\end{ejemplo}

\subsection{Interpretación gráfica}
Si $(a,b)$ pertenece a la gráfica de $f$, entonces $(b,a)$ pertenece a la gráfica de $f^{-1}$. Por ello las gráficas de $y=f(x)$ y $y=f^{-1}(x)$ son simétricas respecto de
\[
\boxed{y=x}.
\]

\subsection{Restricción del dominio}
No toda función es invertible en su dominio original. Para $f(x)=x^2$ sobre $\mathbb R$ se tiene $f(x)=f(-x)$, de modo que no es inyectiva. Si restringimos
\[
f:[0,\infty)\to[0,\infty),\qquad f(x)=x^2,
\]
entonces es biyectiva y
\[
\boxed{f^{-1}(x)=\sqrt{x}.}
\]
No se escribe $\pm\sqrt{x}$: una función debe asignar un único valor a cada entrada.

\begin{ejemplo}
Para
\[
f(x)=(x-2)^2+1,
\]
restringida a $[2,\infty)$, la imagen es $[1,\infty)$. De
\[
y=(x-2)^2+1
\]
se obtiene, usando $x-2\ge0$,
\[
x=2+\sqrt{y-1}.
\]
Así,
\[
\boxed{f^{-1}(x)=2+\sqrt{x-1}.}
\]
\end{ejemplo}

\subsection{Monotonía e inyectividad}
Si una función es estrictamente creciente o estrictamente decreciente en un intervalo, entonces es inyectiva en ese intervalo:
\[
\boxed{\text{estrictamente monótona}\Longrightarrow\text{inyectiva}.}
\]
Esto explica, por ejemplo, por qué $x^3$ es inyectiva en $\mathbb R$ y por qué $x^2$ puede hacerse inyectiva restringiendo su dominio a $[0,\infty)$ o a $(-\infty,0]$.

\subsection{Resumen conceptual}
\[
\boxed{\text{Inyectiva}\iff f(x_1)=f(x_2)\Rightarrow x_1=x_2}
\]
\[
\boxed{\text{Suprayectiva}\iff\operatorname{Im}(f)=B}
\]
\[
\boxed{\text{Biyectiva}\iff\text{inyectiva y suprayectiva}}
\]
\[
\boxed{f:A\to B\text{ posee inversa }B\to A\iff f\text{ es biyectiva}.}
\]

% =========================================================
\section{Banco de problemas: inyectividad, biyectividad e inversas}
% =========================================================

\subsection{Reconocimiento de inyectividad}
Determine si cada función es inyectiva en el dominio indicado. Justifique.
\begin{enumerate}
\item $f:\mathbb R\to\mathbb R$, $f(x)=2x+1$.
\item $f:\mathbb R\to\mathbb R$, $f(x)=x^2$.
\item $f:\mathbb R\to\mathbb R$, $f(x)=x^3$.
\item $f:[0,\infty)\to\mathbb R$, $f(x)=x^2$.
\item $f:\mathbb R\to\mathbb R$, $f(x)=|x|$.
\item $f:[0,\infty)\to[0,\infty)$, $f(x)=\sqrt{x}$.
\item $f:\mathbb R\to(0,\infty)$, $f(x)=e^x$.
\item $f:(0,\infty)\to\mathbb R$, $f(x)=\ln x$.
\end{enumerate}

\subsection{Clasificación: inyectiva, suprayectiva o biyectiva}
Clasifique cada función respecto del dominio y codominio indicados.
\begin{enumerate}
\item $f:\mathbb R\to\mathbb R$, $f(x)=x+4$.
\item $f:\mathbb R\to\mathbb R$, $f(x)=x^2$.
\item $f:\mathbb R\to[0,\infty)$, $f(x)=x^2$.
\item $f:[0,\infty)\to[0,\infty)$, $f(x)=x^2$.
\item $f:\mathbb R\to\mathbb R$, $f(x)=x^3$.
\item $f:\mathbb R\to(0,\infty)$, $f(x)=e^x$.
\item $f:\mathbb R\to\mathbb R$, $f(x)=e^x$.
\item $f:(0,\infty)\to\mathbb R$, $f(x)=\ln x$.
\end{enumerate}

\subsection{Cálculo de funciones inversas}
Determine $f^{-1}$ e indique dominio e imagen de $f$ y de $f^{-1}$.
\begin{enumerate}
\item $f(x)=2x+5$.
\item $f(x)=7-3x$.
\item $f(x)=\dfrac{x-1}{4}$.
\item $f(x)=\dfrac{2x+1}{3}$.
\item $f(x)=x^3$.
\item $f(x)=x^3-4$.
\item $f(x)=\sqrt{x}$, $x\ge0$.
\item $f(x)=e^x$.
\item $f(x)=2^x$.
\item $f(x)=\ln x$.
\end{enumerate}

\subsection{Funciones racionales}
Encuentre la inversa, especificando las restricciones necesarias.
\begin{enumerate}
\item $f(x)=\dfrac1x$.
\item $f(x)=\dfrac{x+1}{x-2}$.
\item $f(x)=\dfrac{2x-3}{x+4}$.
\item $f(x)=\dfrac{3x+1}{2x-5}$.
\item $f(x)=\dfrac{x-2}{x+2}$.
\item $f(x)=1+\dfrac1x$.
\end{enumerate}

\subsection{Restricción del dominio}
En cada caso explique por qué la función no es inyectiva en $\mathbb R$, proponga una restricción adecuada y determine la inversa.
\begin{enumerate}
\item $f(x)=x^2$.
\item $f(x)=(x-1)^2$.
\item $f(x)=(x+3)^2-2$.
\item $f(x)=x^4$.
\item $f(x)=|x|$.
\item $f(x)=\sqrt{x^2}$.
\end{enumerate}

\subsection{Verificación por composición}
Encuentre $f^{-1}$ y compruebe $f(f^{-1}(x))=x$ y $f^{-1}(f(x))=x$.
\begin{enumerate}
\item $f(x)=4x+3$.
\item $f(x)=\dfrac{x-5}{2}$.
\item $f(x)=x^3+2$.
\item $f(x)=e^{2x}$.
\item $f(x)=3^x$.
\item $f(x)=\ln(x+1)$, $x>-1$.
\end{enumerate}

\subsection{Preguntas para discusión}
\begin{enumerate}
\item ¿Puede una función ser inyectiva pero no suprayectiva? Dé un ejemplo.
\item ¿Puede ser suprayectiva pero no inyectiva? Dé un ejemplo.
\item ¿Por qué una función no inyectiva no puede poseer una inversa funcional sobre toda su imagen sin restringir el dominio?
\item ¿Por qué $f^{-1}(x)$ no significa $1/f(x)$?
\item Si $f(a)=b$, ¿qué debe cumplir $f^{-1}$?
\item Explique geométricamente la simetría de $f$ y $f^{-1}$ respecto de $y=x$.
\end{enumerate}


% =========================================================
\section*{Bibliografía básica}
\addcontentsline{toc}{section}{Bibliografía básica}

\begin{itemize}
    \item Larson, Hostetler y Edwards, \emph{Cálculo I}, 7a ed.
    \item Leithold, \emph{El Cálculo}, 7a ed.
    \item Stewart, \emph{Cálculo. Trascendentes tempranas}, 4a ed.
    \item Thomas y Finney, \emph{Cálculo. Una variable}, 9a ed.
    \item Swokowski, \emph{Cálculo con Geometría Analítica}, 2a ed.
\end{itemize}

\end{document}
